\documentclass[12pt, a4paper, oneside]{report}

\usepackage[utf8]{inputenc}
\usepackage[T1]{fontenc}
\usepackage[italian]{babel}
\usepackage[margin=2.5cm]{geometry}
\usepackage{parskip}
\usepackage{graphicx}
\graphicspath{{./immagini/}}
\usepackage{hyperref}
\usepackage{bookmark}
\hypersetup{colorlinks=true, linkcolor=black, urlcolor=blue}
\usepackage{tabularx}
\usepackage{xcolor}
\usepackage[dvipsnames]{xcolor}
\usepackage{ltablex}
\keepXColumns 
\setcounter{secnumdepth}{0}
\usepackage{titlesec}
\titleformat{\chapter}[block]
  {\normalfont\Huge\bfseries} 
  {} 
  {0pt} 
  {} 
\titlespacing*{\chapter}{0pt}{-20pt}{20pt}

\begin{document}

\begin{titlepage}
    \centering
    \null \vfill
    {\huge\bfseries Progetto di Ingegneria del Software\newline Parte A\par}
    {\Large Raffaele Principe (749438) -- Riccardo Marini (746230)\par}
    {\large Anno Accademico 2025/2026\par}
    \vfill
\end{titlepage}

\tableofcontents 
\clearpage

\chapter{Versione 1}

\section{Casi d'uso testuali}

\renewcommand{\arraystretch}{1.4} 

\begin{tabularx}{\textwidth}{|l|X|}
    \hline
    \textbf{Nome}               & \textbf{UC01 - Avvio e Login} \\ \hline
    \textbf{Attore}             & Configuratore \\ \hline
    \endfirsthead 
    \hline
    \multicolumn{2}{|l|}{\textit{Continua dalla pagina precedente...}} \\ \hline
    \endhead 
    \hline
    \multicolumn{2}{|r|}{\textit{Continua nella pagina successiva...}} \\ \hline
    \endfoot 
    \hline
    \endlastfoot
    \textbf{Scenario principale} & 
    Precondizione: L'applicazione è avviata. \newline
    Trigger: L'utente seleziona "Configuratore" dal menu iniziale.\newline
    1. Il sistema mostra la schermata di login dedicata.\newline
    2. Il configuratore inserisce username e password.\newline
    3. Il sistema verifica la corrispondenza delle credenziali.\newline
    4. Il sistema autentica il configuratore.\newline
    5. Il sistema verifica se si tratta del primissimo accesso assoluto (credenziali di default).\newline
    6. <<include>> "UC02 - Imposta credenziali personali" (se primo accesso).\newline
    7. Il sistema mostra il menu principale.\newline
    Postcondizione: Il configuratore è autenticato nel back-end. \newline
    Fine \\ \hline
    \textbf{Scenario alternativo} & 
    2a. Precondizione: il configuratore digita "esci" come username.\newline
    Il sistema interrompe l'operazione e torna al menu di selezione ruolo. \newline
    Fine \\ \hline
    \textbf{Scenario alternativo} & 
    3a. Precondizione: le credenziali non sono valide. \newline
    Il sistema mostra un errore e indica i tentativi rimasti. Torna al punto 2 (se tentativi $<$ 3), altrimenti il sistema blocca l'accesso alla sessione. \newline
    Fine \\ \hline
    \textbf{Scenario alternativo} & 
    3b. Precondizione: nessun configuratore registrato nel sistema.\newline
    Il sistema accetta le credenziali di default (admin/admin123), crea il profilo iniziale, imposta la necessità di cambio credenziali e prosegue al punto 6. \\ \hline
\end{tabularx}

\renewcommand{\arraystretch}{1.4} 

\begin{tabularx}{\textwidth}{|l|X|}
    \hline
    \textbf{Nome}               & \textbf{UC02 - Imposta credenziali personali} \\ \hline
    \textbf{Attore}             & Configuratore \\ \hline
    \endfirsthead 
    \hline
    \multicolumn{2}{|l|}{\textit{Continua dalla pagina precedente...}} \\ \hline
    \endhead 

    \hline
    \multicolumn{2}{|r|}{\textit{Continua nella pagina successiva...}} \\ \hline
    \endfoot 

    \hline
    \endlastfoot

    \textbf{Scenario principale} & 
    Precondizione: Login effettuato con credenziali di default (primo accesso).\newline
    Trigger: Il sistema rileva la necessità di cambiare le credenziali.\newline
    1. Il sistema richiede di scegliere le credenziali personali.\newline
    2. Il configuratore inserisce il nuovo username.\newline
    3. Il sistema verifica che lo username non coincida con la parola riservata "esci".\newline
    4. Il configuratore inserisce la nuova password.\newline
    5. Il configuratore inserisce la password per conferma.\newline
    6. Il sistema verifica che le due password coincidano.\newline
    7. Il sistema verifica l'unicità globale dello username tra Configuratori.\newline
    8. Il sistema aggiorna le credenziali in memoria.\newline
    9. Il sistema salva i dati in modo persistente.\newline
    10. Il sistema mostra la conferma e accede al menu.\newline
    Postcondizione: Le credenziali di default sono sostituite e salvate in modo permanente.\newline
    Fine \\ \hline
    
    \textbf{Scenario alternativo} & 
    2a. Precondizione: il configuratore digita "annulla" nel campo username.\newline
    Il sistema annulla l'operazione ed esegue il logout automatico.\newline
    Fine \\ \hline

    \textbf{Scenario alternativo} & 
    3a. Precondizione: lo username inserito è "esci".\newline
    Il sistema mostra un errore, annulla l'operazione ed esegue il logout.\newline
    Fine \\ \hline

    \textbf{Scenario alternativo} & 
    4a. Precondizione: il configuratore digita "annulla" nel campo password.\newline
    Il sistema annulla l'operazione ed esegue il logout automatico.\newline
    Fine \\ \hline

    \textbf{Scenario alternativo} & 
    6a. Precondizione: le password non coincidono.\newline
    Il sistema mostra un errore, annulla l'operazione ed esegue il logout.\newline
    Fine \\ \hline

    \textbf{Scenario alternativo} & 
    7a. Precondizione: campi lasciati vuoti.\newline
    Il sistema mostra un errore, annulla l'operazione ed esegue il logout.\newline
    Fine \\ \hline

    \textbf{Scenario alternativo} & 
    8a. Precondizione: username già in uso nel sistema.\newline
    Il sistema mostra un errore, annulla l'operazione ed esegue il logout.\newline
    Fine \\ \hline

    \textbf{Scenario alternativo} & 
    9a. Precondizione: Errore di salvataggio (I/O).\newline
    Il sistema ripristina le credenziali in memoria, mostra un errore ed esegue il logout.\newline
    Fine \\ \hline
\end{tabularx}

\renewcommand{\arraystretch}{1.4} 

\begin{tabularx}{\textwidth}{|l|X|}
    \hline
    \textbf{Nome}               & \textbf{UC03 - Logout} \\ \hline
    \textbf{Attore}             & Configuratore \\ \hline
    \endfirsthead 

    \hline
    \multicolumn{2}{|l|}{\textit{Continua dalla pagina precedente...}} \\ \hline
    \endhead 

    \hline
    \multicolumn{2}{|r|}{\textit{Continua nella pagina successiva...}} \\ \hline
    \endfoot 

    \hline
    \endlastfoot

    \textbf{Scenario principale} & 
    Precondizione: Configuratore autenticato. \newline
    Trigger: Selezione "Logout". \newline
    1. Il configuratore seleziona Logout. \newline
    2. Il sistema termina la sessione corrente. \newline
    3. Il sistema chiede se accedere con un altro account. \newline
    4. Sì: ritorna a UC01. \newline
    5. No: termina l'applicazione. \newline
    Fine \\ \hline
    
    \textbf{Scenario alternativo} & 
    3a. Precondizione: Il configuratore inserisce una risposta non valida (diversa da "s" o "n").\newline
    Il sistema mostra un messaggio di errore testuale e ripropone la domanda. \\ \hline
\end{tabularx}

\renewcommand{\arraystretch}{1.4} 

\begin{tabularx}{\textwidth}{|l|X|}
    \hline
    \textbf{Nome}               & \textbf{UC04 - Visualizza campi base} \\ \hline
    \textbf{Attore}             & Configuratore \\ \hline
    \endfirsthead 

    \hline
    \multicolumn{2}{|l|}{\textit{Continua dalla pagina precedente...}} \\ \hline
    \endhead 

    \hline
    \multicolumn{2}{|r|}{\textit{Continua nella pagina successiva...}} \\ \hline
    \endfoot 

    \hline
    \endlastfoot

    \textbf{Scenario principale} & 
    Precondizione: Il configuratore è autenticato. I campi base sono stati inizializzati.\newline
    Trigger: Il configuratore seleziona "Gestione campi base" dal menu principale.\newline
    1. Il configuratore seleziona la voce "Gestione campi base".\newline
    2. Il sistema mostra l'elenco dei campi base con nome e tipo, evidenziando che sono tutti obbligatori (Titolo, Numero partecipanti, ecc.).\newline
    3. Il sistema indica che i campi base sono immutabili.\newline
    4. Il configuratore preme INVIO per tornare al menu principale.\newline
    Fine \\ \hline
    
    \textbf{Scenario alternativo} & 
    2a. Precondizione: Campi base non inizializzati.\newline
    Il sistema li inizializza automaticamente al primo avvio e poi mostra l'elenco.\newline
    Fine \\ \hline
\end{tabularx}

\renewcommand{\arraystretch}{1.4} 

\begin{tabularx}{\textwidth}{|l|X|}
    \hline
    \textbf{Nome}               & \textbf{UC05 - Aggiungi campo comune} \\ \hline
    \textbf{Attore}             & Configuratore \\ \hline
    \endfirsthead 

    \hline
    \multicolumn{2}{|l|}{\textit{Continua dalla pagina precedente...}} \\ \hline
    \endhead 

    \hline
    \multicolumn{2}{|r|}{\textit{Continua nella pagina successiva...}} \\ \hline
    \endfoot 

    \hline
    \endlastfoot

    \textbf{Scenario principale} & 
    Precondizione: Il configuratore è autenticato. \newline
    Trigger: Il configuratore seleziona "Aggiungi campo comune" nel menu "Gestione campi comuni". \newline
    1. Il configuratore seleziona "Aggiungi campo comune". \newline
    2. Il sistema chiede il nome del nuovo campo. \newline
    3. Il configuratore inserisce il nome. \newline
    4. Il sistema chiede se il campo è obbligatorio o facoltativo. \newline
    5. Il configuratore risponde (s/n). \newline
    6. Il sistema verifica che il nome non coincida con un campo base o comune esistente. \newline
    7. Il sistema aggiunge il campo comune e salva in modo persistente. \newline
    8. Il sistema mostra conferma dell'operazione. \newline
    Fine \\ \hline
    
    \textbf{Scenario alternativo} & 
    3a. Precondizione: Nome vuoto. \newline
    Il sistema mostra errore e richiede nuovamente il nome. \\ \hline

    \textbf{Scenario alternativo} & 
    6a. Precondizione: Conflitto con campo base. \newline
    Il sistema mostra errore e non aggiunge il campo. \newline
    Fine \\ \hline

    \textbf{Scenario alternativo} & 
    6b. Precondizione: Conflitto con campo comune esistente. \newline
    Il sistema mostra errore e non aggiunge il campo. \newline
    Fine \\ \hline

    \textbf{Scenario alternativo} & 
    7a. Precondizione: Errore di salvataggio (I/O). \newline
    Il sistema aggiunge il campo in memoria ma avvisa dell'errore di salvataggio fisico. \newline
    Fine \\ \hline
\end{tabularx}

\renewcommand{\arraystretch}{1.4} 

\begin{tabularx}{\textwidth}{|l|X|}
    \hline
    \textbf{Nome}               & \textbf{UC06 - Rimuovi campo comune} \\ \hline
    \textbf{Attore}             & Configuratore \\ \hline
    \endfirsthead 

    \hline
    \multicolumn{2}{|l|}{\textit{Continua dalla pagina precedente...}} \\ \hline
    \endhead 

    \hline
    \multicolumn{2}{|r|}{\textit{Continua nella pagina successiva...}} \\ \hline
    \endfoot 

    \hline
    \endlastfoot

    \textbf{Scenario principale} & 
    Precondizione: Il configuratore è autenticato. Esiste almeno un campo comune. \newline
    Trigger: Il configuratore seleziona "Rimuovi campo comune". \newline
    1. Il configuratore seleziona "Rimuovi campo comune". \newline
    2. Il sistema mostra l'elenco dei campi comuni esistenti. \newline
    3. Il configuratore inserisce il nome del campo da rimuovere. \newline
    4. Il sistema rimuove il campo comune e salva in modo persistente. \newline
    5. Il sistema mostra conferma dell'operazione. \newline
    Fine \\ \hline
    
    \textbf{Scenario alternativo} & 
    3a. Precondizione: Campo non trovato. \newline
    Il sistema mostra errore e non esegue la rimozione. \newline
    Fine \\ \hline

    \textbf{Scenario alternativo} & 
    4a. Precondizione: Errore di salvataggio (I/O). \newline
    Il sistema rimuove il campo in memoria ma avvisa dell'errore di salvataggio fisico. \newline
    Fine \\ \hline
\end{tabularx}

\renewcommand{\arraystretch}{1.4} 

\begin{tabularx}{\textwidth}{|l|X|}
    \hline
    \textbf{Nome}               & \textbf{UC07 - Modifica obbligatorietà campo comune} \\ \hline
    \textbf{Attore}             & Configuratore \\ \hline
    \endfirsthead 

    \hline
    \multicolumn{2}{|l|}{\textit{Continua dalla pagina precedente...}} \\ \hline
    \endhead 

    \hline
    \multicolumn{2}{|r|}{\textit{Continua nella pagina successiva...}} \\ \hline
    \endfoot 

    \hline
    \endlastfoot

    \textbf{Scenario principale} & 
    Precondizione: Il configuratore è autenticato. Esiste almeno un campo comune. \newline
    Trigger: Il configuratore seleziona "Modifica obbligatorietà campo comune". \newline
    1. Il configuratore seleziona "Modifica obbligatorietà campo comune". \newline
    2. Il sistema mostra l'elenco dei campi comuni con l'obbligatorietà attuale. \newline
    3. Il configuratore inserisce il nome del campo da modificare. \newline
    4. Il configuratore indica il nuovo valore di obbligatorietà (s/n). \newline
    5. Il sistema aggiorna il campo e salva in modo persistente. \newline
    6. Il sistema mostra conferma dell'operazione. \newline
    Fine \\ \hline
    
    \textbf{Scenario alternativo} & 
    3a. Precondizione: Campo non trovato. \newline
    Il sistema mostra errore e non applica modifiche. \newline
    Fine \\ \hline

    \textbf{Scenario alternativo} & 
    5a. Precondizione: Errore di salvataggio (I/O). \newline
    Il sistema aggiorna il campo in memoria ma avvisa dell'errore di salvataggio fisico. \newline
    Fine \\ \hline
\end{tabularx}

\renewcommand{\arraystretch}{1.4} 

\begin{tabularx}{\textwidth}{|l|X|}
    \hline
    \textbf{Nome}               & \textbf{UC08 - Aggiungi categoria} \\ \hline
    \textbf{Attore}             & Configuratore \\ \hline
    \endfirsthead 

    \hline
    \multicolumn{2}{|l|}{\textit{Continua dalla pagina precedente...}} \\ \hline
    \endhead 

    \hline
    \multicolumn{2}{|r|}{\textit{Continua nella pagina successiva...}} \\ \hline
    \endfoot 

    \hline
    \endlastfoot

    \textbf{Scenario principale} & 
    Precondizione: Il configuratore è autenticato. I campi base sono inizializzati. \newline
    Trigger: Il configuratore seleziona "Aggiungi categoria". \newline
    1. Il configuratore seleziona "Aggiungi categoria". \newline
    2. Il sistema chiede il nome della nuova categoria. \newline
    3. Il configuratore inserisce il nome. \newline
    4. Il sistema verifica che il nome non sia già in uso. \newline
    5. Il sistema crea la categoria e salva in modo persistente. \newline
    6. Il sistema chiede se aggiungere subito campi specifici. \newline
    7. Sì: il caso d'uso estende UC10 (aggiunta campi specifici). \newline
    8. No: il sistema torna al menu categorie. \newline
    9. Il sistema mostra conferma dell'operazione. \newline
    Fine \\ \hline
    
    \textbf{Scenario alternativo} & 
    3a. Precondizione: Nome vuoto. \newline
    Errore e richiesta nuovo inserimento. \\ \hline

    \textbf{Scenario alternativo} & 
    4a. Precondizione: Categoria duplicata. \newline
    Errore e nessuna creazione. \newline
    Fine \\ \hline

    \textbf{Scenario alternativo} & 
    5a. Precondizione: Errore di salvataggio (I/O). \newline
    Il sistema crea la categoria in memoria ma avvisa dell'errore di salvataggio fisico. \newline
    Fine \\ \hline
\end{tabularx}

\renewcommand{\arraystretch}{1.4} 

\begin{tabularx}{\textwidth}{|l|X|}
    \hline
    \textbf{Nome}               & \textbf{UC09 - Rimuovi categoria} \\ \hline
    \textbf{Attore}             & Configuratore \\ \hline
    \endfirsthead 

    \hline
    \multicolumn{2}{|l|}{\textit{Continua dalla pagina precedente...}} \\ \hline
    \endhead 

    \hline
    \multicolumn{2}{|r|}{\textit{Continua nella pagina successiva...}} \\ \hline
    \endfoot 

    \hline
    \endlastfoot

    \textbf{Scenario principale} & 
    Precondizione: Il configuratore è autenticato. Esiste almeno una categoria. \newline
    Trigger: Il configuratore seleziona "Rimuovi categoria". \newline
    1. Il configuratore seleziona "Rimuovi categoria". \newline
    2. Il sistema mostra l'elenco categorie. \newline
    3. Il configuratore inserisce il nome della categoria da rimuovere. \newline
    4. Il sistema chiede conferma esplicita prima della rimozione. \newline
    5. Il configuratore conferma. \newline
    6. Il sistema rimuove la categoria e tutti i suoi campi specifici. \newline
    7. Il sistema salva i dati in modo persistente. \newline
    8. Il sistema mostra conferma. \newline
    Fine \\ \hline
    
    \textbf{Scenario alternativo} & 
    3a. Precondizione: Categoria non trovata. \newline
    Errore, nessuna rimozione. \newline
    Fine \\ \hline

    \textbf{Scenario alternativo} & 
    5a. Precondizione: Annulla. \newline
    Il configuratore non conferma; operazione interrotta. \newline
    Fine \\ \hline

    \textbf{Scenario alternativo} & 
    7a. Precondizione: Errore di salvataggio (I/O). \newline
    Il sistema rimuove la categoria in memoria ma avvisa dell'errore di salvataggio fisico. \newline
    Fine \\ \hline
\end{tabularx}

\renewcommand{\arraystretch}{1.4} 

\begin{tabularx}{\textwidth}{|l|X|}
    \hline
    \textbf{Nome}               & \textbf{UC10 - Aggiungi campo specifico a categoria} \\ \hline
    \textbf{Attore}             & Configuratore \\ \hline
    \endfirsthead 

    \hline
    \multicolumn{2}{|l|}{\textit{Continua dalla pagina precedente...}} \\ \hline
    \endhead 

    \hline
    \multicolumn{2}{|r|}{\textit{Continua nella pagina successiva...}} \\ \hline
    \endfoot 

    \hline
    \endlastfoot

    \textbf{Scenario principale} & 
    Precondizione: Il configuratore è autenticato. Esiste almeno una categoria. \newline
    Trigger: Il configuratore sceglie di aggiungere un campo specifico a una categoria. \newline
    1. Il configuratore seleziona "Gestisci campi specifici" e indica la categoria. \newline
    2. Il sistema verifica che la categoria esista. \newline
    3. Il configuratore seleziona "Aggiungi campo specifico". \newline
    4. Il sistema chiede il nome del nuovo campo. \newline
    5. Il configuratore inserisce il nome. \newline
    6. Il sistema chiede se il campo è obbligatorio o facoltativo. \newline
    7. Il configuratore risponde (s/n). \newline
    8. Il sistema verifica che il nome non coincida con un campo base, comune o specifico già presente. \newline
    9. Il sistema aggiunge il campo specifico alla categoria e salva. \newline
    10. Il sistema mostra conferma e permette inserimenti successivi (INVIO per terminare).\newline
    Fine \\ \hline
    
    \textbf{Scenario alternativo} & 
    2a. Precondizione: Categoria non esiste. \newline
    Il sistema mostra un messaggio di errore. \newline
    Fine \\ \hline

    \textbf{Scenario alternativo} & 
    5a. Precondizione: Nome vuoto. \newline
    Il sistema mostra un errore e richiede un nuovo inserimento. \\ \hline

    \textbf{Scenario alternativo} & 
    8a. Precondizione: Nome duplicato. \newline
    Il sistema mostra un errore e non procede con l'aggiunta. \newline
    Fine \\ \hline

    \textbf{Scenario alternativo} & 
    9a. Precondizione: Errore di salvataggio (I/O). \newline
    Il sistema aggiunge il campo in memoria ma avvisa dell'errore di salvataggio fisico. \newline
    Fine \\ \hline
\end{tabularx}

\renewcommand{\arraystretch}{1.4} 

\begin{tabularx}{\textwidth}{|l|X|}
    \hline
    \textbf{Nome}               & \textbf{UC11 - Rimuovi campo specifico da categoria} \\ \hline
    \textbf{Attore}             & Configuratore \\ \hline
    \endfirsthead 

    \hline
    \multicolumn{2}{|l|}{\textit{Continua dalla pagina precedente...}} \\ \hline
    \endhead 

    \hline
    \multicolumn{2}{|r|}{\textit{Continua nella pagina successiva...}} \\ \hline
    \endfoot 

    \hline
    \endlastfoot

    \textbf{Scenario principale} & 
    Precondizione: Il configuratore è autenticato. Esiste una categoria con almeno un campo specifico. \newline
    Trigger: Il configuratore seleziona "Rimuovi campo specifico". \newline
    1. Il configuratore seleziona la categoria. \newline
    2. Il sistema mostra i campi specifici della categoria. \newline
    3. Il configuratore inserisce il nome del campo da rimuovere. \newline
    4. Il sistema rimuove il campo specifico e salva in modo persistente. \newline
    5. Il sistema mostra conferma. \newline
    Fine \\ \hline
    
    \textbf{Scenario alternativo} & 
    3a. Precondizione: Campo non trovato. \newline
    Il sistema mostra un messaggio di errore. \newline
    Fine \\ \hline

    \textbf{Scenario alternativo} & 
    4a. Precondizione: Errore di salvataggio (I/O). \newline
    Il sistema rimuove il campo in memoria ma avvisa dell'errore di salvataggio fisico. \newline
    Fine \\ \hline
\end{tabularx}

\renewcommand{\arraystretch}{1.4} 

\begin{tabularx}{\textwidth}{|l|X|}
    \hline
    \textbf{Nome}               & \textbf{UC12 - Modifica obbligatorietà campo specifico} \\ \hline
    \textbf{Attore}             & Configuratore \\ \hline
    \endfirsthead 

    \hline
    \multicolumn{2}{|l|}{\textit{Continua dalla pagina precedente...}} \\ \hline
    \endhead 

    \hline
    \multicolumn{2}{|r|}{\textit{Continua nella pagina successiva...}} \\ \hline
    \endfoot 

    \hline
    \endlastfoot

    \textbf{Scenario principale} & 
    Precondizione: Il configuratore è autenticato. Esiste una categoria con almeno un campo specifico. \newline
    Trigger: Il configuratore seleziona "Modifica obbligatorietà campo specifico". \newline
    1. Il configuratore seleziona la categoria. \newline
    2. Il sistema mostra i campi specifici con obbligatorietà attuale. \newline
    3. Il configuratore inserisce il nome del campo da modificare. \newline
    4. Il configuratore indica il nuovo valore (s/n). \newline
    5. Il sistema aggiorna il campo specifico e salva in modo persistente. \newline
    6. Il sistema mostra conferma. \newline
    Fine \\ \hline
    
    \textbf{Scenario alternativo} & 
    3a. Precondizione: Campo non trovato. \newline
    Il sistema mostra un messaggio di errore. \newline
    Fine \\ \hline

    \textbf{Scenario alternativo} & 
    5a. Precondizione: Errore di salvataggio (I/O). \newline
    Il sistema aggiorna il campo in memoria ma avvisa dell'errore di salvataggio fisico. \newline
    Fine \\ \hline
\end{tabularx}

\renewcommand{\arraystretch}{1.4} 

\begin{tabularx}{\textwidth}{|l|X|}
    \hline
    \textbf{Nome}               & \textbf{UC13 - Visualizza categorie e campi} \\ \hline
    \textbf{Attore}             & Configuratore \\ \hline
    \endfirsthead 

    \hline
    \multicolumn{2}{|l|}{\textit{Continua dalla pagina precedente...}} \\ \hline
    \endhead 

    \hline
    \multicolumn{2}{|r|}{\textit{Continua nella pagina successiva...}} \\ \hline
    \endfoot 

    \hline
    \endlastfoot

    \textbf{Scenario principale} & 
    Precondizione: Il configuratore è autenticato. \newline
    Trigger: Il configuratore seleziona "Visualizza riepilogo categorie e campi". \newline
    1. Il configuratore seleziona "Visualizza riepilogo". \newline
    2. Il sistema mostra i campi BASE. \newline
    3. Il sistema mostra i campi COMUNI (con indicazione di obbligatorio/facoltativo). \newline
    4. Per ogni categoria, il sistema mostra: nome categoria e campi specifici (con indicazione di obbligatorio/facoltativo). \newline
    5. Il configuratore preme INVIO per tornare al menu. \newline
    Fine \\ \hline
    
    \textbf{Scenario alternativo} & 
    3a. Precondizione: Nessun campo comune definito. \newline
    Il sistema mostra la dicitura "(nessuno definito)". \\ \hline

    \textbf{Scenario alternativo} & 
    4a. Precondizione: Nessuna categoria definita nel sistema. \newline
    Il sistema mostra la dicitura "(nessuna categoria definita)". \\ \hline

    \textbf{Scenario alternativo} & 
    4b. Precondizione: Una categoria non possiede campi specifici. \newline
    Il sistema mostra la dicitura "(nessun campo specifico)" sotto il nome della categoria interessata. \\ \hline
\end{tabularx}

\renewcommand{\arraystretch}{1.4} 

\begin{tabularx}{\textwidth}{|l|X|}
    \hline
    \textbf{Nome}               & \textbf{UC14 - Visualizza dettaglio categoria} \\ \hline
    \textbf{Attore}             & Configuratore \\ \hline
    \endfirsthead 

    \hline
    \multicolumn{2}{|l|}{\textit{Continua dalla pagina precedente...}} \\ \hline
    \endhead 

    \hline
    \multicolumn{2}{|r|}{\textit{Continua nella pagina successiva...}} \\ \hline
    \endfoot 

    \hline
    \endlastfoot

    \textbf{Scenario principale} & 
    Precondizione: Il configuratore è autenticato. Esiste almeno una categoria. \newline
    Trigger: Il configuratore seleziona "Visualizza dettaglio categoria". \newline
    1. Il configuratore seleziona "Visualizza dettaglio categoria". \newline
    2. Il sistema richiede l'inserimento del nome della categoria. \newline
    3. Il configuratore inserisce il nome richiesto. \newline
    4. Il sistema mostra, per la categoria selezionata, il set completo dei dati: tutti i campi BASE, i campi COMUNI e i campi SPECIFICI della categoria stessa. \newline
    5. Il configuratore preme INVIO per terminare la visualizzazione e tornare al menu. \newline
    Fine \\ \hline
    
    \textbf{Scenario alternativo} & 
    3a. Precondizione: La categoria inserita non è presente nell'archivio. \newline
    Il sistema mostra un messaggio di errore e l'operazione si interrompe. \newline
    Fine \\ \hline
\end{tabularx}


\section{Casi d'uso grafici}

\begin{figure}[htbp]
    \centering
    \makebox[\textwidth][c]{
        \includegraphics[width=0.8\textwidth]{immagini/casi uso grafici v1 definitivo 2.png} 
    }
    \label{fig:usecase1_v1}
\end{figure}

\clearpage

\begin{figure}[htbp]
    \centering
    \makebox[\textwidth][c]{
        \includegraphics[width=1.3\textwidth]{immagini/casi uso grafici v1 definitivo 1a.png} 
    }
    \label{fig:usecase2_v1}
\end{figure}



\section{Diagramma delle classi}

\begin{figure}[p]
    \centering
    \includegraphics[width=0.9\textwidth]{immagini/UML_V1.png}
    \label{fig:diagramma_classi_v1}
\end{figure}



\chapter{Versione 2}

\section{Casi d'uso testuali modificati}

\renewcommand{\arraystretch}{1.4} 

\begin{tabularx}{\textwidth}{|l|X|}
    \hline
    \textbf{Nome}               & \textbf{UC01 - Avvio e Login} \\ \hline
    \textbf{Attore}             & Configuratore \\ \hline
    \endfirsthead 

    \hline
    \multicolumn{2}{|l|}{\textit{Continua dalla pagina precedente...}} \\ \hline
    \endhead 

    \hline
    \multicolumn{2}{|r|}{\textit{Continua nella pagina successiva...}} \\ \hline
    \endfoot 

    \hline
    \endlastfoot

    \textbf{Scenario principale} & 
    Precondizione: L'applicazione è avviata. \newline
    Trigger: L'utente seleziona "Configuratore" dal menu iniziale.\newline
    1. Il sistema mostra la schermata di login dedicata.\newline
    2. Il configuratore inserisce username e password.\newline
    3. Il sistema verifica la corrispondenza delle credenziali.\newline
    4. Il sistema autentica il configuratore.\newline
    5. Il sistema verifica se si tratta del primissimo accesso assoluto (credenziali di default).\newline
    6. <<include>> "UC02 - Imposta credenziali personali" (se primo accesso).\newline
    \textcolor{BurntOrange}{7. Il sistema mostra il menu principale con i contatori aggiornati: numero di proposte in sessione e numero di proposte aperte in bacheca.}\newline
    Postcondizione: Il configuratore è autenticato nel back-end. \newline
    Fine \\ \hline
    
    \textbf{Scenario alternativo} & 
    2a. Precondizione: il configuratore digita "esci" come username.\newline
    Il sistema interrompe l'operazione e torna al menu di selezione ruolo. \newline
    Fine \\ \hline

    \textbf{Scenario alternativo} & 
    3a. Precondizione: le credenziali non sono valide. \newline
    Il sistema mostra un errore e indica i tentativi rimasti. Torna al punto 2 (se tentativi $<$ 3), altrimenti il sistema blocca l'accesso alla sessione. \newline
    Fine \\ \hline

    \textbf{Scenario alternativo} & 
    3b. Precondizione: nessun configuratore registrato nel sistema.\newline
    Il sistema accetta le credenziali di default (admin/admin123), crea il profilo iniziale, imposta la necessità di cambio credenziali e prosegue al punto 6. \\ \hline
\end{tabularx}

\renewcommand{\arraystretch}{1.4} 

\begin{tabularx}{\textwidth}{|l|X|}
    \hline
    \textbf{Nome}               & \textbf{UC03 - Logout} \\ \hline
    \textbf{Attore}             & Configuratore \\ \hline
    \endfirsthead 

    \hline
    \multicolumn{2}{|l|}{\textit{Continua dalla pagina precedente...}} \\ \hline
    \endhead 

    \hline
    \multicolumn{2}{|r|}{\textit{Continua nella pagina successiva...}} \\ \hline
    \endfoot 

    \hline
    \endlastfoot

    \textbf{Scenario principale} & 
    Precondizione: Configuratore autenticato. \newline
    Trigger: Selezione "Logout". \newline
    1. Il configuratore seleziona Logout. \newline
    \textcolor{ForestGreen}{2. Il sistema scarta le proposte non pubblicate della sessione corrente (stato BOZZA o VALIDA).}\newline
    \textcolor{ForestGreen}{3. Se erano presenti proposte non pubblicate, il sistema avvisa il configuratore.}\newline
    4. Il sistema termina la sessione corrente. \newline
    5. Il sistema chiede se accedere con un altro account. \newline
    6. Sì: ritorna a UC01. \newline
    7. No: termina l'applicazione. \newline
    Fine \\ \hline
    
    \textbf{Scenario alternativo} & 
    5a. Precondizione: Il configuratore inserisce una risposta non valida (diversa da "s" o "n").\newline
    Il sistema mostra un messaggio di errore testuale e ripropone la domanda. \\ \hline
\end{tabularx}

\section{Casi d'uso testuali nuovi}

\renewcommand{\arraystretch}{1.4} 

\begin{tabularx}{\textwidth}{|l|X|}
    \hline
    \textbf{Nome}               & \textbf{UC15 - Crea proposta} \\ \hline
    \textbf{Attore}             & Configuratore \\ \hline
    \endfirsthead 

    \hline
    \multicolumn{2}{|l|}{\textit{Continua dalla pagina precedente...}} \\ \hline
    \endhead 

    \hline
    \multicolumn{2}{|r|}{\textit{Continua nella pagina successiva...}} \\ \hline
    \endfoot 

    \hline
    \endlastfoot

    \textbf{Scenario principale} & 
    Precondizione: Il configuratore è autenticato ed esiste almeno una categoria.\newline
    Trigger: Selezione "Nuova proposta" dal menu gestionale.\newline
    1. Il sistema mostra l'elenco delle categorie disponibili.\newline
    2. Il configuratore inserisce il nome della categoria desiderata.\newline
    3. Il sistema verifica l'esistenza della categoria.\newline
    4. Il sistema cattura lo snapshot dei campi (base, comuni e specifici) rendendoli indipendenti da future modifiche alla configurazione.\newline
    5. Il sistema assegna un ID univoco progressivo estratto dal contatore globale.\newline
    6. Il sistema aggiunge la proposta alla sessione corrente in stato BOZZA.\newline
    7. <<include>> "UC16 - Compila campi proposta".\newline
    8. Il sistema mostra il riepilogo e, se in stato VALIDA, offre la pubblicazione immediata (estende UC17).\newline
    Fine \\ \hline
    
    \textbf{Scenario alternativo} & 
    1a. Precondizione: l'archivio categorie è vuoto.\newline
    Il sistema mostra l'errore "Nessuna categoria disponibile" e impedisce la creazione.\newline
    Fine \\ \hline

    \textbf{Scenario alternativo} & 
    3a. Precondizione: categoria non trovata in base all'input dell'utente.\newline
    Il sistema mostra un errore e torna al menu proposte.\newline
    Fine \\ \hline
\end{tabularx}

\renewcommand{\arraystretch}{1.4} 

\begin{tabularx}{\textwidth}{|l|X|}
    \hline
    \textbf{Nome}               & \textbf{UC16 - Compila / Modifica campi proposta} \\ \hline
    \textbf{Attore}             & Configuratore \\ \hline
    \endfirsthead 

    \hline
    \multicolumn{2}{|l|}{\textit{Continua dalla pagina precedente...}} \\ \hline
    \endhead 

    \hline
    \multicolumn{2}{|r|}{\textit{Continua nella pagina successiva...}} \\ \hline
    \endfoot 

    \hline
    \endlastfoot

    \textbf{Scenario principale} & 
    Precondizione: Esiste almeno una proposta in sessione.\newline
    Trigger: Estensione da UC15 o selezione "Continua compilazione".\newline
    1. Il sistema presenta i campi dello snapshot uno alla volta.\newline
    2. Per ogni campo, il sistema mostra l'etichetta [OBB/FAC], il formato richiesto e il valore attuale.\newline
    3. Il configuratore inserisce un valore o preme INVIO per mantenere il valore attuale.\newline
    4. Il sistema valida l'input e ricalcola automaticamente lo stato (BOZZA/VALIDA).\newline
    5. Il sistema verifica i vincoli temporali: Termine iscrizione $>$ oggi; Data inizio $\ge$ Termine + 2gg; Data fine $\ge$ Data inizio; Formato Ora (HH:MM).\newline
    6. Il sistema mostra il riepilogo finale e l'elenco degli errori se lo stato è BOZZA.\newline
    Fine \\ \hline
    
    \textbf{Scenario alternativo} & 
    4a. Precondizione: formato data (gg/mm/aaaa) o ora (HH:MM) errato.\newline
    Il sistema mostra un errore specifico e non aggiorna il valore del campo. Prosegue con il campo successivo. \\ \hline
\end{tabularx}

\renewcommand{\arraystretch}{1.4} 

\begin{tabularx}{\textwidth}{|l|X|}
    \hline
    \textbf{Nome}               & \textbf{UC17 - Pubblica proposta in bacheca} \\ \hline
    \textbf{Attore}             & Configuratore \\ \hline
    \endfirsthead 

    \hline
    \multicolumn{2}{|l|}{\textit{Continua dalla pagina precedente...}} \\ \hline
    \endhead 

    \hline
    \multicolumn{2}{|r|}{\textit{Continua nella pagina successiva...}} \\ \hline
    \endfoot 

    \hline
    \endlastfoot

    \textbf{Scenario principale} & 
    Precondizione: Il configuratore è autenticato e c'è almeno una proposta VALIDA in sessione.\newline
    Trigger: Selezione "Pubblica proposta in bacheca".\newline
    1. Il sistema mostra l'elenco delle proposte in sessione attualmente VALIDE.\newline
    2. Il configuratore inserisce l'ID della proposta da pubblicare.\newline
    3. Il sistema cattura la data odierna.\newline
    4. Il sistema imposta lo stato della proposta su APERTA e registra la data nello storico.\newline
    5. Il sistema sposta la proposta dalla sessione all'Archivio del configuratore.\newline
    6. Il sistema salva le modifiche in modo persistente.\newline
    7. Il sistema mostra conferma dell'operazione.\newline
    Postcondizione: La proposta ha stato APERTA, è salvata nell'archivio ed è visibile ai fruitori. \newline
    Fine \\ \hline
    
    \textbf{Scenario alternativo} & 
    1a. Precondizione: nessuna proposta in sessione è in stato VALIDA. \newline
    Il sistema mostra un messaggio di errore e l'operazione termina.\newline
    Fine \\ \hline

    \textbf{Scenario alternativo} & 
    2a. Precondizione: l'input non è numerico o non corrisponde a un ID presente in elenco. \newline
    Il sistema mostra un errore e l'operazione termina.\newline
    Fine \\ \hline
\end{tabularx}

\renewcommand{\arraystretch}{1.4} 

\begin{tabularx}{\textwidth}{|l|X|}
    \hline
    \textbf{Nome}               & \textbf{UC18 - Elimina proposta dalla sessione} \\ \hline
    \textbf{Attore}             & Configuratore \\ \hline
    \endfirsthead 

    \hline
    \multicolumn{2}{|l|}{\textit{Continua dalla pagina precedente...}} \\ \hline
    \endhead 

    \hline
    \multicolumn{2}{|r|}{\textit{Continua nella pagina successiva...}} \\ \hline
    \endfoot 

    \hline
    \endlastfoot

    \textbf{Scenario principale} & 
    Precondizione: Il configuratore è autenticato. Esiste almeno una proposta in sessione (stato BOZZA o VALIDA).\newline
    Trigger: Il configuratore seleziona "Elimina proposta" nel menu proposte.\newline
    1. Il sistema mostra l'elenco delle proposte in sessione.\newline
    2. Il configuratore inserisce l'ID della proposta da eliminare.\newline
    3. Il sistema chiede conferma esplicita prima dell'eliminazione.\newline
    4. Il configuratore conferma.\newline
    5. Il sistema rimuove definitivamente la proposta dalla sessione.\newline
    6. Il sistema mostra conferma dell'avvenuta operazione.\newline
    Fine \\ \hline
    
    \textbf{Scenario alternativo} & 
    2a. Precondizione: l'ID inserito non è valido o non corrisponde a una proposta presente in sessione.\newline
    Il sistema mostra un messaggio di errore e interrompe l'operazione.\newline
    Fine \\ \hline

    \textbf{Scenario alternativo} & 
    4a. Precondizione: il configuratore annulla l'operazione (nega la conferma).\newline
    Il sistema interrompe l'operazione senza apportare modifiche e torna al menu.\newline
    Fine \\ \hline
\end{tabularx}

\renewcommand{\arraystretch}{1.4} 

\begin{tabularx}{\textwidth}{|l|X|}
    \hline
    \textbf{Nome}               & \textbf{UC19 - Visualizza bacheca} \\ \hline
    \textbf{Attore}             & Configuratore \\ \hline
    \endfirsthead 

    \hline
    \multicolumn{2}{|l|}{\textit{Continua dalla pagina precedente...}} \\ \hline
    \endhead 

    \hline
    \multicolumn{2}{|r|}{\textit{Continua nella pagina successiva...}} \\ \hline
    \endfoot 

    \hline
    \endlastfoot

    \textbf{Scenario principale} & 
    Precondizione: Il configuratore è autenticato.\newline
    Trigger: Selezione "Visualizza bacheca".\newline
    1. Il configuratore richiede la visualizzazione della bacheca.\newline
    2. Il sistema recupera tutte le proposte con stato APERTA.\newline
    3. Il sistema raggruppa le proposte per categoria esistente.\newline
    4. Il sistema mostra per ogni proposta: ID, Data pubblicazione, Creatore e i valori di tutti i campi compilati (omettendo i campi facoltativi lasciati vuoti).\newline
    5. Il configuratore preme INVIO per terminare la visualizzazione.\newline
    Postcondizione: Nessuna modifica allo stato del sistema. \newline
    Fine \\ \hline
    
    \textbf{Scenario alternativo} & 
    2a. Precondizione: non ci sono proposte in stato APERTA nell'archivio.\newline
    Il sistema mostra il messaggio "(Nessuna proposta aperta in bacheca)".\newline
    Fine \\ \hline

    \textbf{Scenario alternativo} & 
    3a. Precondizione: la categoria originaria di una proposta è stata rimossa dopo la pubblicazione.\newline
    Il sistema visualizza la proposta raggruppandola sotto l'etichetta "(categoria rimossa)".\newline
    Fine \\ \hline
\end{tabularx}

\section{Casi d'uso grafici}
\clearpage
\begin{figure}[htbp]
    \centering
    \makebox[\textwidth][c]{
        \includegraphics[width=1.3\textwidth]{immagini/casi uso grafici v2 definitivo.png} 
    }
    \label{fig:usecase_v2}
\end{figure}
\clearpage
\section{Diagramma delle classi}
\clearpage
\begin{figure}[htbp]
    \centering
    \makebox[\textwidth][c]{
        \includegraphics[width=0.84\textwidth]{immagini/diagramma_v2.png}
    }
    \label{fig:diagramma_classi_v2}
\end{figure}


\chapter{Versione 3}

\section{Casi d'uso testuali modificati}

\begin{table}[htbp]
    \centering
    \renewcommand{\arraystretch}{1.4} 
    \begin{tabularx}{\textwidth}{|l|X|}
    \hline
    \textbf{Nome}               & \textbf{UC01 - Avvio e Login} \\ \hline
    \textbf{Attore}             & Configuratore \\ \hline
    \textbf{Scenario principale} & 
    Precondizione: L'applicazione è avviata. \textcolor{ForestGreen}{Transizioni automatiche (UC20) completate dal sistema}.\newline
    Trigger: L'utente seleziona "Configuratore" dal menu iniziale.\newline
    1. Il sistema mostra la schermata di login dedicata.\newline
    2. Il configuratore inserisce username e password.\newline
    3. Il sistema verifica la corrispondenza delle credenziali.\newline
    4. Il sistema autentica il configuratore.\newline
    5. Il sistema verifica se si tratta del primissimo accesso assoluto (credenziali di default).\newline
    6. <<include>> "UC02 - Imposta credenziali personali" (se primo accesso).\newline
    \textcolor{BurntOrange}{7. Il sistema mostra il menu principale aggiornato coi contatori in tempo reale}.\newline
    Postcondizione: Il configuratore è autenticato nel back-end.\newline
    Fine \\ \hline
    \textbf{Scenario alternativo} & 
    2a. Precondizione: il configuratore digita "esci" come username.\newline
    Il sistema interrompe l'operazione e torna al menu di selezione ruolo.\newline
    Fine \\ \hline
    \textbf{Scenario alternativo} & 
    3a. Precondizione: le credenziali non sono valide.\newline
    Il sistema mostra un errore e indica i tentativi rimasti.\newline
    Torna al punto 2 (se tentativi $< 3$), altrimenti il sistema blocca l'accesso alla sessione.\newline
    Fine \\ \hline
    \end{tabularx}
    \label{tab:uc01_v3}
\end{table}

\renewcommand{\arraystretch}{1.4} 

\begin{tabularx}{\textwidth}{|l|X|}
    \hline
    \textbf{Nome}               & \textbf{UC02 - Imposta credenziali personali} \\ \hline
    \textbf{Attore}             & Configuratore \\ \hline
    \endfirsthead 
    \hline
    \multicolumn{2}{|l|}{\textit{Continua dalla pagina precedente...}} \\ \hline
    \endhead 
    \hline
    \multicolumn{2}{|r|}{\textit{Continua nella pagina successiva...}} \\ \hline
    \endfoot 
    \hline
    \endlastfoot 
    \textbf{Scenario principale} & 
    Precondizione: Login effettuato con credenziali di default (primo accesso).\newline
    Trigger: Il sistema rileva la necessità di cambiare le credenziali.\newline
    1. Il sistema richiede di scegliere le credenziali personali.\newline
    2. Il configuratore inserisce il nuovo username.\newline
    3. Il sistema verifica che lo username non coincida con la parola riservata "esci".\newline
    4. Il configuratore inserisce la nuova password.\newline
    5. Il configuratore inserisce la password per conferma.\newline
    6. Il sistema verifica che le due password coincidano.\newline
    \textcolor{ForestGreen}{7. Il sistema verifica che username e password non siano stringhe vuote/blank.}\newline
    \textcolor{BurntOrange}{8. Il sistema verifica l'unicità globale dello username (tra Configuratori e Fruitori).}\newline
    9. Il sistema aggiorna le credenziali in memoria.\newline
    10. Il sistema salva i dati in modo persistente.\newline
    11. Il sistema mostra la conferma e accede al menu.\newline
    Postcondizione: Le credenziali di default sono sostituite e salvate in modo permanente.\newline
    Fine \\ \hline
    \textbf{Scenario alternativo} & 
    2a. Precondizione: il configuratore digita "annulla" nel campo username.\newline
    Il sistema annulla l'operazione ed esegue il logout automatico.\newline
    Fine \\ \hline
    \textbf{Scenario alternativo} & 
    3a. Precondizione: lo username inserito è "esci".\newline
    Il sistema mostra un errore, annulla l'operazione ed esegue il logout.\newline
    Fine \\ \hline
    \textbf{Scenario alternativo} & 
    4a. Precondizione: il configuratore digita "annulla" nel campo password.\newline
    Il sistema annulla l'operazione ed esegue il logout automatico.\newline
    Fine \\ \hline
    \textbf{Scenario alternativo} & 
    6a. Precondizione: le password non coincidono.\newline
    Il sistema mostra un errore, annulla l'operazione ed esegue il logout.\newline
    Fine \\ \hline
    \textbf{Scenario alternativo} & 
    7a. Precondizione: campi lasciati vuoti.\newline
    Il sistema mostra un errore, annulla l'operazione ed esegue il logout.\newline
    Fine \\ \hline
    \textbf{Scenario alternativo} & 
    8a. Precondizione: username già in uso nel sistema.\newline
    Il sistema mostra un errore, annulla l'operazione ed esegue il logout.\newline
    Fine \\ \hline
    \textbf{Scenario alternativo} & 
    \textcolor{ForestGreen}{10a. Precondizione: Errore di salvataggio (I/O).}\newline
    \textcolor{ForestGreen}{Il sistema esegue un rollback in memoria, mostra un errore ed esegue il logout.}\newline
    \textcolor{ForestGreen}{Fine} \\ \hline
\end{tabularx}

\renewcommand{\arraystretch}{1.4} 

\begin{tabularx}{\textwidth}{|l|X|}
    \hline
    \textbf{Nome}               & \textbf{UC17 - Pubblica proposta in bacheca} \\ \hline
    \textbf{Attore}             & Configuratore \\ \hline
    \endfirsthead 
    \hline
    \multicolumn{2}{|l|}{\textit{Continua dalla pagina precedente...}} \\ \hline
    \endhead 
    \hline
    \multicolumn{2}{|r|}{\textit{Continua nella pagina successiva...}} \\ \hline
    \endfoot 
    \hline
    \endlastfoot
    \textbf{Scenario principale} & 
    Precondizione: Il configuratore è autenticato e c'è almeno una proposta VALIDA in sessione.\newline
    Trigger: Selezione "Pubblica proposta in bacheca".\newline
    1. Il sistema mostra l'elenco delle proposte in sessione attualmente VALIDE.\newline
    2. Il configuratore inserisce l'ID della proposta da pubblicare.\newline
    3. Il sistema cattura la data odierna.\newline
    \textcolor{ForestGreen}{4. Il sistema ricalcola la validità della proposta in base alla data odierna.}\newline
    5. Il sistema imposta lo stato su APERTA e registra la data nello storico.\newline
    \textcolor{BurntOrange}{6. Il sistema sposta la proposta dalla sessione all'Archivio globale.}\newline
    7. Il sistema salva le modifiche in modo persistente.\newline
    8. Il sistema mostra conferma dell'operazione.\newline
    Postcondizione: La proposta ha stato APERTA, è salvata nell'archivio ed è visibile ai fruitori.\newline
    Fine \\ \hline
    \textbf{Scenario alternativo} & 
    1a. Precondizione: nessuna proposta in sessione è in stato VALIDA.\newline
    Il sistema mostra un errore.\newline
    Fine \\ \hline
    \textbf{Scenario alternativo} & 
    2a. Precondizione: l'input non è numerico o non corrisponde a una proposta in elenco.\newline
    Il sistema mostra un errore.\newline
    Fine \\ \hline
    \textbf{Scenario alternativo} & 
    \textcolor{ForestGreen}{4a. Precondizione: la validità è scaduta a causa del cambio di data (es. superata mezzanotte).}\newline
    \textcolor{ForestGreen}{Il sistema nega la pubblicazione e mostra l'elenco degli errori.}\newline
    \textcolor{ForestGreen}{Fine} \\ \hline
    \textbf{Scenario alternativo} & 
    \textcolor{ForestGreen}{7a. Precondizione: Errore di salvataggio (I/O).}\newline
    \textcolor{ForestGreen}{Il sistema esegue un rollback completo riportando la proposta allo stato VALIDA in sessione.}\newline
    \textcolor{ForestGreen}{Fine} \\ \hline
\end{tabularx}

\renewcommand{\arraystretch}{1.4} 

\begin{tabularx}{\textwidth}{|l|X|}
    \hline
    \textbf{Nome}               & \textbf{UC19 - Visualizza bacheca Configuratore} \\ \hline
    \textbf{Attore}             & Configuratore \\ \hline
    \endfirsthead 
    \hline
    \multicolumn{2}{|l|}{\textit{Continua dalla pagina precedente...}} \\ \hline
    \endhead 
    \hline
    \multicolumn{2}{|r|}{\textit{Continua nella pagina successiva...}} \\ \hline
    \endfoot 
    \hline
    \endlastfoot
    \textbf{Scenario principale} & 
    Precondizione: Il configuratore è autenticato.\newline
    Trigger: Selezione "Visualizza bacheca".\newline
    1. Il configuratore richiede la visualizzazione della bacheca.\newline
    2. Il sistema recupera tutte le proposte con stato APERTA.\newline
    \textcolor{BurntOrange}{3. Il sistema raggruppa dinamicamente le proposte per categoria (inclusa la voce "(categoria rimossa)" per le eventuali proposte orfane).}\newline
    \textcolor{BurntOrange}{4. Il sistema mostra per ogni proposta: ID, dettagli, valori dei campi e il numero di iscritti attuali rispetto al massimo consentito.}\newline
    5. Il configuratore preme INVIO per terminare.\newline
    Postcondizione: Nessuna modifica allo stato del sistema.\newline
    Fine \\ \hline
    \textbf{Scenario alternativo} & 
    2a. Precondizione: non ci sono proposte APERTE in archivio.\newline
    Il sistema mostra il messaggio "(Nessuna proposta aperta in bacheca)".\newline
    Fine \\ \hline
    \textbf{Scenario alternativo} & 
    3a. Precondizione: categoria di una proposta rimossa dopo la pubblicazione.\newline
    Il sistema visualizza la proposta sotto l'etichetta "(categoria rimossa)".\newline
    Fine \\ \hline
\end{tabularx}


\section{Casi d'uso testuali nuovi}

\renewcommand{\arraystretch}{1.4} 

\begin{tabularx}{\textwidth}{|l|X|}
    \hline
    \textbf{Nome}               & \textbf{UC20 - Esecuzione transizioni automatiche di stato} \\ \hline
    \textbf{Attore}             & Sistema \\ \hline
    \endfirsthead 
    \hline
    \multicolumn{2}{|l|}{\textit{Continua dalla pagina precedente...}} \\ \hline
    \endhead 
    \hline
    \multicolumn{2}{|r|}{\textit{Continua nella pagina successiva...}} \\ \hline
    \endfoot 
    \hline
    \endlastfoot
    \textbf{Scenario principale} & 
    Precondizione: L'applicazione viene avviata. L'archivio contiene proposte.\newline
    Trigger: Inizializzazione del sistema.\newline
    1. Il sistema analizza l'archivio delle proposte.\newline
    2. Per ogni proposta APERTA con "Termine ultimo" superato, il sistema verifica se iscritti $\ge$ minimo richiesto.\newline
    3. Il sistema transita la proposta in stato CONFERMATA.\newline
    4. Il sistema genera una notifica di conferma per ogni aderente iscritto.\newline
    5. Per ogni proposta CONFERMATA con "Data conclusiva" superata, il sistema la transita in stato CONCLUSA.\newline
    6. Il sistema salva i nuovi stati e le notifiche in modo persistente.\newline
    7. Il sistema mostra a schermo il numero totale di transizioni effettuate.\newline
    Postcondizione: Stati aggiornati alla data odierna e notifiche distribuite.\newline
    Fine \\ \hline
    \textbf{Scenario alternativo} & 
    2a. Precondizione: la proposta APERTA è scaduta ma gli iscritti sono insufficienti (o il campo "Numero massimo" è corrotto/negativo).\newline
    Il sistema transita la proposta in ANNULLATA e genera una notifica di annullamento per ogni aderente.\newline
    Torna al punto 5 \\ \hline
    \textbf{Scenario alternativo} & 
    6a. Precondizione: Errore di salvataggio (I/O).\newline
    Il sistema esegue un rollback ricaricando i dati precedenti dal disco per evitare la generazione di notifiche duplicate ai riavvii successivi.\newline
    Fine \\ \hline
\end{tabularx}

\renewcommand{\arraystretch}{1.4} 

\begin{tabularx}{\textwidth}{|l|X|}
    \hline
    \textbf{Nome}               & \textbf{UC21 - Registrazione Fruitore} \\ \hline
    \textbf{Attore}             & Fruitore \\ \hline
    \endfirsthead 
    \hline
    \multicolumn{2}{|l|}{\textit{Continua dalla pagina precedente...}} \\ \hline
    \endhead 
    \hline
    \multicolumn{2}{|r|}{\textit{Continua nella pagina successiva...}} \\ \hline
    \endfoot 
    \hline
    \endlastfoot
    \textbf{Scenario principale} & 
    Precondizione: L'utente non è autenticato.\newline
    Trigger: Selezione "Registrazione (nuovo account)" dal menu di accesso.\newline
    1. Il sistema mostra la schermata di registrazione.\newline
    2. Il fruitore inserisce lo username desiderato.\newline
    3. Il fruitore inserisce la password.\newline
    4. Il fruitore reinserisce la password per conferma.\newline
    5. Il sistema verifica che le due password coincidano.\newline
    6. Il sistema verifica che username e password non siano vuoti.\newline
    7. Il sistema verifica l'unicità globale dello username.\newline
    8. Il sistema crea il profilo Fruitore in memoria.\newline
    9. Il sistema salva i dati in modo persistente.\newline
    10. Il sistema autentica automaticamente il fruitore e mostra il benvenuto.\newline
    Postcondizione: Il fruitore è registrato, autenticato e con uno spazio notifiche vuoto.\newline
    Fine \\ \hline
    \textbf{Scenario alternativo} & 
    2a. Precondizione: il fruitore digita "annulla" come username.\newline
    Il sistema interrompe l'operazione e torna al menu.\newline
    Fine \\ \hline
    \textbf{Scenario alternativo} & 
    3a. Precondizione: il fruitore digita "annulla" come password.\newline
    Il sistema interrompe l'operazione e torna al menu.\newline
    Fine \\ \hline
    \textbf{Scenario alternativo} & 
    5a. Precondizione: le password non coincidono.\newline
    Il sistema mostra l'errore e interrompe l'operazione.\newline
    Fine \\ \hline
    \textbf{Scenario alternativo} & 
    6a. Precondizione: campi vuoti o composti da soli spazi.\newline
    Il sistema mostra l'errore e interrompe l'operazione.\newline
    Fine \\ \hline
    \textbf{Scenario alternativo} & 
    7a. Precondizione: username già in uso (globale).\newline
    Il sistema mostra l'errore e interrompe l'operazione.\newline
    Fine \\ \hline 
    \textbf{Scenario alternativo} & 
    9a. Precondizione: Errore di salvataggio (I/O).\newline
    Il sistema rimuove il profilo dalla memoria (rollback), mostra l'errore e interrompe l'operazione.\newline
    Fine \\ \hline
\end{tabularx}

\renewcommand{\arraystretch}{1.4} 

\begin{tabularx}{\textwidth}{|l|X|}
    \hline
    \textbf{Nome}               & \textbf{UC22 - Login Fruitore} \\ \hline
    \textbf{Attore}             & Fruitore \\ \hline
    \endfirsthead 
    \hline
    \multicolumn{2}{|l|}{\textit{Continua dalla pagina precedente...}} \\ \hline
    \endhead 
    \hline
    \multicolumn{2}{|r|}{\textit{Continua nella pagina successiva...}} \\ \hline
    \endfoot 
    \hline
    \endlastfoot
    \textbf{Scenario principale} & 
    Precondizione: Il fruitore possiede un account e non è autenticato.\newline
    Trigger: Selezione "Login (account esistente)".\newline
    1. Il sistema mostra la schermata di login.\newline
    2. Il fruitore inserisce username e password.\newline
    3. Il sistema verifica la corrispondenza delle credenziali.\newline
    4. Il sistema autentica il fruitore.\newline
    5. Il sistema conteggia le notifiche presenti nello spazio personale e mostra il benvenuto con il relativo avviso.\newline
    Postcondizione: Il fruitore è autenticato nel sistema.\newline
    Fine \\ \hline
    \textbf{Scenario alternativo} & 
    2a. Precondizione: il fruitore digita "annulla" nel campo username.\newline
    Il sistema annulla l'operazione e torna al menu.\newline
    Fine \\ \hline
    \textbf{Scenario alternativo} & 
    3a. Precondizione: le credenziali sono errate (incluso inserire "annulla" come password).\newline
    Il sistema mostra l'errore e i tentativi rimasti.\newline
    Torna al punto 2 (se tentativi $<$ 3), altrimenti il sistema blocca l'accesso.\newline
    Fine \\ \hline
\end{tabularx}

\renewcommand{\arraystretch}{1.4} 

\begin{tabularx}{\textwidth}{|l|X|}
    \hline
    \textbf{Nome}               & \textbf{UC23 - Visualizza Bacheca (Fruitore)} \\ \hline
    \textbf{Attore}             & Fruitore \\ \hline
    \endfirsthead 
    \hline
    \multicolumn{2}{|l|}{\textit{Continua dalla pagina precedente...}} \\ \hline
    \endhead 
    \hline
    \multicolumn{2}{|r|}{\textit{Continua nella pagina successiva...}} \\ \hline
    \endfoot 
    \hline
    \endlastfoot
    \textbf{Scenario principale} & 
    Precondizione: Il fruitore è autenticato.\newline
    Trigger: Selezione "Visualizza bacheca".\newline
    1. Il fruitore richiede la visualizzazione.\newline
    2. Il sistema recupera tutte le proposte in stato APERTA.\newline
    3. Il sistema raggruppa le proposte per categoria.\newline
    4. Il sistema mostra dettagli, posti occupati, scadenze e tutti i campi informativi per ogni proposta.\newline
    5. Il sistema evidenzia con l'etichetta \textcolor{red}{[ISCRITTO]} le iniziative a cui l'utente loggato ha già aderito.\newline
    6. Il fruitore preme INVIO per terminare la consultazione.\newline
    Postcondizione: Nessuna modifica allo stato del sistema.\newline
    Fine \\ \hline  
    \textbf{Scenario alternativo} & 
    2a. Precondizione: nessuna proposta APERTA trovata nell'archivio.\newline
    Il sistema comunica "(Nessuna proposta aperta in bacheca)".\newline
    Fine \\ \hline
\end{tabularx}

\renewcommand{\arraystretch}{1.4} 

\begin{tabularx}{\textwidth}{|l|X|}
    \hline
    \textbf{Nome}               & \textbf{UC24 - Iscrizione a proposta} \\ \hline
    \textbf{Attore}             & Fruitore \\ \hline
    \endfirsthead 
    \hline
    \multicolumn{2}{|l|}{\textit{Continua dalla pagina precedente...}} \\ \hline
    \endhead 
    \hline
    \multicolumn{2}{|r|}{\textit{Continua nella pagina successiva...}} \\ \hline
    \endfoot 
    \hline
    \endlastfoot
    \textbf{Scenario principale} & 
    Precondizione: Il fruitore è autenticato.\newline
    Trigger: Selezione "Aderisci a una proposta".\newline
    1. Il fruitore avvia la procedura di iscrizione.\newline
    2. Il sistema estrae le proposte APERTE dall'archivio.\newline
    3. Il sistema verifica che ci sia almeno una proposta APERTA nella bacheca globale.\newline
    4. Il sistema filtra le proposte escludendo quelle a cui l'utente è già iscritto.\newline
    5. Il sistema verifica che la lista filtrata non sia vuota.\newline
    6. Il fruitore inserisce l'ID della proposta scelta.\newline
    7. Il sistema verifica che l'input sia numerico e diverso da "0".\newline
    8. Il sistema individua la proposta e verifica che sia in stato APERTA.\newline
    9. Il sistema verifica che la data odierna sia $\le$ al "Termine ultimo di iscrizione".\newline
    10. Il sistema verifica che ci siano posti disponibili.\newline
    11. Il sistema aggiunge lo username del fruitore alla lista degli aderenti.\newline
    12. Il sistema salva i dati in modo persistente.\newline
    13. Il sistema conferma l'avvenuta iscrizione all'utente.\newline
    Postcondizione: L'utente è iscritto e i posti disponibili scalano.\newline
    Fine \\ \hline
    \textbf{Scenario alternativo} & 
    3a. Precondizione: l'archivio globale non contiene alcuna proposta in stato APERTA.\newline
    Il sistema lo comunica all'utente e termina l'operazione.\newline
    Fine \\ \hline
    \textbf{Scenario alternativo} & 
    5a. Precondizione: il fruitore è già iscritto a tutte le proposte APERTE (lista filtrata vuota).\newline
    Il sistema lo comunica all'utente e termina l'operazione.\newline
    Fine \\ \hline
    \textbf{Scenario alternativo} & 
    7a. Precondizione: il fruitore digita "0" per annullare l'inserimento.\newline
    Il sistema annulla l'operazione.\newline
    Fine \\ \hline
    \textbf{Scenario alternativo} & 
    7b. Precondizione: l'input inserito non è un numero intero.\newline
    Il sistema mostra il messaggio "ID non valido".\newline
    Fine \\ \hline
    \textbf{Scenario alternativo} & 
    8a. Precondizione: l'ID non corrisponde a nessuna proposta o la proposta non è più APERTA.\newline
    Il sistema mostra un messaggio di errore.\newline
    Fine \\ \hline
    \textbf{Scenario alternativo} & 
    9a. Precondizione: il termine di iscrizione è scaduto rispetto alla data odierna.\newline
    Il sistema blocca l'iscrizione e mostra un messaggio di errore.\newline
    Fine \\ \hline
    \textbf{Scenario alternativo} & 
    10a. Precondizione: i posti sono esauriti (o il campo numero partecipanti è corrotto).\newline
    Il sistema blocca l'iscrizione e mostra un messaggio di errore.\newline
    Fine \\ \hline
    \textbf{Scenario alternativo} & 
    12a. Precondizione: Errore di salvataggio (I/O).\newline
    Il sistema annulla l'iscrizione in memoria (rollback) e mostra un messaggio di errore.\newline
    Fine \\ \hline
\end{tabularx}

\renewcommand{\arraystretch}{1.4} 

\begin{tabularx}{\textwidth}{|l|X|}
    \hline
    \textbf{Nome}               & \textbf{UC25 - Gestione Spazio Personale (Notifiche)} \\ \hline
    \textbf{Attore}             & Fruitore \\ \hline
    \endfirsthead 
    \hline
    \multicolumn{2}{|l|}{\textit{Continua dalla pagina precedente...}} \\ \hline
    \endhead 
    \hline
    \multicolumn{2}{|r|}{\textit{Continua nella pagina successiva...}} \\ \hline
    \endfoot 
    \hline
    \endlastfoot
    \textbf{Scenario principale} & 
    Precondizione: Il fruitore è autenticato.\newline
    Trigger: Selezione "Spazio personale".\newline
    1. Il fruitore accede allo spazio personale.\newline
    2. Il sistema mostra l'elenco delle notifiche e le opzioni del menu ("c", "t", "0").\newline
    3. Il fruitore inserisce l'opzione "c" (Cancella singola notifica).\newline
    4. Il sistema verifica che la lista delle notifiche non sia vuota.\newline
    5. Il sistema richiede l'inserimento dell'ID della notifica.\newline
    6. Il fruitore inserisce l'ID.\newline
    7. Il sistema individua e rimuove la notifica dalla memoria.\newline
    8. Il sistema salva i dati in modo persistente.\newline
    9. Il sistema conferma la cancellazione all'utente.\newline
    Torna al punto 2 \\ \hline
    \textbf{Scenario alternativo} & 
    2a. Precondizione: la lista delle notifiche dell'utente è vuota.\newline
    Il sistema stampa "(Nessuna notifica)" ma mostra comunque il menu operativo.\newline
    Prosegue al punto 3 \\ \hline
    \textbf{Scenario alternativo} & 
    3a. Precondizione: il fruitore seleziona "t" (Cancellazione totale).\newline
    Il sistema verifica che la lista non sia vuota, chiede conferma (s/n). L'utente risponde "s". Il sistema svuota la lista in memoria.\newline
    Salta al punto 8 \\ \hline
    \textbf{Scenario alternativo} & 
    3b. Precondizione: il fruitore seleziona "t", ma annulla la conferma (risponde "n").\newline
    Il sistema annulla l'operazione.\newline
    Torna al punto 2 \\ \hline
    \textbf{Scenario alternativo} & 
    3c. Precondizione: il fruitore seleziona "0. Torna al menu principale".\newline
    Il sistema chiude lo spazio personale e torna al menu precedente.\newline
    Fine \\ \hline
    \textbf{Scenario alternativo} & 
    3d. Precondizione: il fruitore inserisce un'opzione non valida (es. lettera errata).\newline
    Il sistema mostra l'errore "Scelta non valida".\newline
    Torna al punto 2 \\ \hline
    \textbf{Scenario alternativo} & 
    4a. Precondizione: il fruitore seleziona "c" o "t", ma la lista delle notifiche è vuota.\newline
    Il sistema mostra l'errore "Nessuna notifica da cancellare".\newline
    Torna al punto 2 \\ \hline
    \textbf{Scenario alternativo} & 
    7a. Precondizione: l'ID inserito non è numerico o non corrisponde ad alcuna notifica in memoria.\newline
    Il sistema mostra un errore.\newline
    Torna al punto 2 \\ \hline
    \textbf{Scenario alternativo} & 
    8a. Precondizione: Errore di salvataggio (I/O).\newline
    Il sistema ripristina in memoria la lista originale intatta (rollback) e mostra un errore.\newline
    Torna al punto 2 \\ \hline
\end{tabularx}

\renewcommand{\arraystretch}{1.4} 

\begin{tabularx}{\textwidth}{|l|X|}
    \hline
    \textbf{Nome}               & \textbf{UC26 - Visualizza Archivio Proposte} \\ \hline
    \textbf{Attore}             & Configuratore \\ \hline
    \endfirsthead 
    \hline
    \multicolumn{2}{|l|}{\textit{Continua dalla pagina precedente...}} \\ \hline
    \endhead 
    \hline
    \multicolumn{2}{|r|}{\textit{Continua nella pagina successiva...}} \\ \hline
    \endfoot 
    \hline
    \endlastfoot
    \textbf{Scenario principale} & 
    Precondizione: Il configuratore è autenticato.\newline
    Trigger: Selezione "Visualizza archivio proposte".\newline
    1. Il configuratore avvia la richiesta di visualizzazione.\newline
    2. Il sistema recupera l'intero archivio (escludendo gli stati interni BOZZA e VALIDA).\newline
    3. Il sistema raggruppa dinamicamente le proposte in base allo stato (APERTA, CONFERMATA, ANNULLATA, CONCLUSA, RITIRATA).\newline
    4. Il sistema mostra i dettagli anagrafici, i posti occupati, la lista degli aderenti e lo storico completo dei passaggi di stato per ogni proposta.\newline
    5. Il configuratore preme INVIO per terminare.\newline
    Postcondizione: Nessuna modifica allo stato del sistema.\newline
    Fine \\ \hline 
    \textbf{Scenario alternativo} & 
    2a. Precondizione: l'archivio è vuoto.\newline
    Il sistema segnala "(Nessuna proposta nell'archivio)" e termina l'operazione.\newline
    Fine \\ \hline
\end{tabularx}


\section{Casi d'uso grafici}

\clearpage
\begin{figure}[htbp]
    \centering
    \makebox[\textwidth][c]{
        \includegraphics[width=1.3\textwidth]{immagini/casi uso grafici v3 definitivo 1.png} 
    }
    \label{fig:usecase_v31}
\end{figure}

\clearpage
\begin{figure}[htbp]
    \centering
    \makebox[\textwidth][c]{
        \includegraphics[width=1.3\textwidth]{immagini/casi uso grafici v3 definitivo 2.png} 
    }
    \label{fig:usecase_v32}
\end{figure}
\section{Diagramma delle classi}
\clearpage
\begin{figure}[p]
    \centering
    \includegraphics[width=\textwidth]{immagini/UML V3 View-Controller-Persistence.png}
    \label{fig:diagramma_classi_v3_view_controller_persistence}
\end{figure}

\begin{figure}[p]
    \centering
    \includegraphics[width=0.9\textwidth]{immagini/UML V3 Model.png}
    \label{fig:diagramma_classi_v3_model}
\end{figure}


% --- CAPITOLO 4 ---
\chapter{Versione 4}
\section{Casi d'uso testuali nuovi}

\renewcommand{\arraystretch}{1.4} 

\begin{tabularx}{\textwidth}{|l|X|}
    \hline
    \textbf{Nome}               & \textbf{UC27 - Ritira proposta in archivio} \\ \hline
    \textbf{Attore}             & Configuratore \\ \hline
    \endfirsthead 

    \hline
    \multicolumn{2}{|l|}{\textit{Continua dalla pagina precedente...}} \\ \hline
    \endhead 

    \hline
    \multicolumn{2}{|r|}{\textit{Continua nella pagina successiva...}} \\ \hline
    \endfoot 

    \hline
    \endlastfoot

    \textbf{Scenario principale} & 
    Precondizione: Il configuratore è autenticato. L'archivio contiene almeno una proposta in stato APERTA o CONFERMATA.\newline
    Trigger: Selezione dell'opzione "Ritira proposta".\newline
    1. Il sistema mostra le proposte in stato APERTA o CONFERMATA presenti nell'archivio.\newline
    2. Il configuratore inserisce l'ID della proposta da ritirare.\newline
    3. Il sistema estrae la proposta e verifica che il suo stato attuale consenta il ritiro.\newline
    4. Il sistema verifica che la data odierna sia strettamente precedente alla "Data inizio" dell'evento.\newline
    5. Il sistema transita lo stato della proposta in RITIRATA e registra il passaggio nello storico.\newline
    6. Il sistema invia automaticamente una notifica di ritiro nello spazio personale di tutti i fruitori iscritti (aderenti).\newline
    7. Il sistema salva le modifiche in modo persistente.\newline
    8. Il sistema mostra la conferma dell'operazione. \newline
    Fine \\ \hline
    
    \textbf{Scenario alternativo} & 
    2a. Precondizione: l'ID inserito non corrisponde ad alcuna proposta o la proposta non è nell'archivio.\newline
    Il sistema mostra un messaggio di errore e interrompe l'operazione. \newline
    Fine \\ \hline

    \textbf{Scenario alternativo} & 
    3a. Precondizione: la proposta si trova in uno stato diverso da APERTA o CONFERMATA (es. già CONCLUSA o ANNULLATA).\newline
    Il sistema mostra l'errore "Il ritiro è consentito solo per proposte APERTE o CONFERMATE" e interrompe l'operazione. \newline
    Fine \\ \hline

    \textbf{Scenario alternativo} & 
    4a. Precondizione: la "Data inizio" dell'evento è già trascorsa o corrisponde alla data odierna (oggi $\ge$ Data inizio).\newline
    Il sistema blocca il ritiro indicando che l'evento è imminente o passato. \newline
    Fine \\ \hline

    \textbf{Scenario alternativo} & 
    7a. Precondizione: Errore di salvataggio (I/O).\newline
    Il sistema esegue un rollback ricaricando i dati precedenti dal disco per ripristinare sia lo stato della proposta che le notifiche, poi mostra un errore. \newline
    Fine \\ \hline
\end{tabularx}

\renewcommand{\arraystretch}{1.4} 

\begin{tabularx}{\textwidth}{|l|X|}
    \hline
    \textbf{Nome}               & \textbf{UC28 - Disdici iscrizione a proposta} \\ \hline
    \textbf{Attore}             & Fruitore \\ \hline
    \endfirsthead 

    \hline
    \multicolumn{2}{|l|}{\textit{Continua dalla pagina precedente...}} \\ \hline
    \endhead 

    \hline
    \multicolumn{2}{|r|}{\textit{Continua nella pagina successiva...}} \\ \hline
    \endfoot 

    \hline
    \endlastfoot

    \textbf{Scenario principale} & 
    Precondizione: Il fruitore è autenticato e risulta iscritto ad almeno una proposta.\newline
    Trigger: Selezione dell'opzione per disdire un'iscrizione.\newline
    1. Il sistema mostra l'elenco delle proposte APERTE a cui l'utente è iscritto.\newline
    2. Il fruitore inserisce l'ID della proposta da cui desidera disiscriversi.\newline
    3. Il sistema verifica che la proposta sia in stato APERTA.\newline
    4. Il sistema verifica che la data odierna sia minore o uguale al "Termine ultimo di iscrizione".\newline
    5. Il sistema rimuove lo username del fruitore dalla lista degli aderenti della proposta.\newline
    6. Il sistema salva le modifiche in modo persistente.\newline
    7. Il sistema mostra la conferma dell'avvenuta disdetta. \newline
    Fine \\ \hline
    
    \textbf{Scenario alternativo} & 
    2a. Precondizione: l'ID non è valido oppure il fruitore non risulta iscritto alla proposta indicata.\newline
    Il sistema mostra un errore. \newline
    Fine \\ \hline

    \textbf{Scenario alternativo} & 
    3a. Precondizione: la proposta ha cambiato stato (non è più APERTA).\newline
    Il sistema mostra un errore indicando l'impossibilità di operare. \newline
    Fine \\ \hline

    \textbf{Scenario alternativo} & 
    4a. Precondizione: il "Termine ultimo di iscrizione" è scaduto rispetto alla data odierna.\newline
    Il sistema blocca la disdetta informando che il termine è oltrepassato. \newline
    Fine \\ \hline

    \textbf{Scenario alternativo} & 
    6a. Precondizione: Errore di salvataggio (I/O).\newline
    Il sistema esegue un rollback reinserendo il fruitore tra gli aderenti in memoria e mostra un messaggio di errore. \newline
    Fine \\ \hline
\end{tabularx}

\section{Casi d'uso grafici}


\begin{figure}[htbp]
    \centering
    \makebox[\textwidth][c]{
        \includegraphics[width=1.3\textwidth]{immagini/casi uso grafici v4 definitivo 1.png} 
    }
    \label{fig:usecase_v41}
\end{figure}


\begin{figure}[htbp]
    \centering
    \makebox[\textwidth][c]{
        \includegraphics[width=1.3\textwidth]{immagini/casi uso grafici v4 definitivo 2.png} 
    }
    \label{fig:usecase_v42}
\end{figure}

\clearpage
\section{Diagramma delle classi}
\begin{figure}[p]
    \centering
    \includegraphics[width=0.8\textwidth]{immagini/uml v4 2.png}
    \label{fig:diagramma_classi_v4_view_controller_persistence}
\end{figure}

\begin{figure}[p]
    \centering
    \includegraphics[width=1.14\textwidth]{immagini/uml V4 model.png}
    \label{fig:diagramma_classi_v4_model}
\end{figure}
% --- CAPITOLO 5 ---
\chapter{Versione 5}
\section{Casi d'uso testuali nuovi}
\renewcommand{\arraystretch}{1.4} 

\begin{tabularx}{\textwidth}{|l|X|}
    \hline
    \textbf{Nome} & \textbf{UC29 – Importa singolo file batch} \\ \hline
    \textbf{Attore} & Configuratore \\ \hline
    \endfirsthead 

    \hline
    \multicolumn{2}{|l|}{\textit{Continua dalla pagina precedente...}} \\ \hline
    \endhead 

    \hline
    \multicolumn{2}{|r|}{\textit{Continua nella pagina successiva...}} \\ \hline
    \endfoot 

    \hline
    \endlastfoot

    \textbf{Scenario principale} & 
    Precondizione: Il configuratore è autenticato. Il file batch esiste nel filesystem ed è leggibile.\newline
    Trigger: Selezione dell'opzione "Importa da file batch" $\rightarrow$ "Importa singolo file".\newline
    1. Il sistema mostra la descrizione del formato batch supportato (CAMPO\_COMUNE, CATEGORIA, PROPOSTA).\newline
    2. Il configuratore inserisce il percorso del file (assoluto o relativo alla working directory).\newline
    3. Il sistema legge tutte le righe del file.\newline
    4. Il sistema scorre le righe, ignorando quelle vuote e quelle che iniziano con \#.\newline
    5. Per ogni riga significativa, il sistema identifica il comando (primo token prima di |).\newline
    6. Per ogni riga CAMPO\_COMUNE: il sistema crea il campo comune con nome e obbligatorietà indicati e salva su disco.\newline
    7. Per ogni riga CATEGORIA: il sistema crea la categoria; se la riga contiene token CAMPO\_SPECIFICO, aggiunge i campi specifici inline prima del salvataggio, poi salva su disco.\newline
    8. Per ogni riga PROPOSTA: il sistema costruisce lo snapshot dei campi (BASE $\rightarrow$ COMUNI $\rightarrow$ SPECIFICI della categoria), assegna i valori indicati con il formato campo = valore, valida la proposta con le stesse regole della modalità interattiva e, se valida, la pubblica direttamente in bacheca salvando su disco.\newline
    9. Al termine, il sistema raccoglie il resoconto in un \textit{BatchRisultato} (numero righe elaborate, successi, warning, errori).\newline
    10. Il sistema mostra il resoconto completo all'utente.\newline
    Fine \\ \hline
    
    \textbf{Scenario alternativo} & 
    2a. Precondizione: il percorso è vuoto o l'utente digita "annulla".\newline
    Il sistema annulla l'operazione senza modifiche.\newline
    Fine \\ \hline

    \textbf{Scenario alternativo} & 
    3a. Precondizione: il file non esiste o non è accessibile (IOException).\newline
    Il sistema mostra un messaggio di errore con indicazione del percorso e interrompe l'operazione.\newline
    Fine \\ \hline

    \textbf{Scenario alternativo} & 
    5a. Precondizione: il comando della riga non è riconosciuto.\newline
    Il sistema registra un warning con numero di riga e prosegue con la riga successiva.\newline
    Fine (riga) \\ \hline

    \textbf{Scenario alternativo} & 
    6a. Precondizione: il campo comune ha nome vuoto, oppure esiste già un campo con lo stesso nome.\newline
    Il sistema registra un warning e salta la riga.\newline
    Fine (riga) \\ \hline

    \textbf{Scenario alternativo} & 
    6b. Precondizione: errore di salvataggio (I/O) dopo l'aggiunta in memoria.\newline
    Il sistema esegue il rollback in memoria, registra un errore e prosegue.\newline
    Fine (riga) \\ \hline

    \textbf{Scenario alternativo} & 
    7a. Precondizione: la categoria esiste già.\newline
    Il sistema registra un warning e salta la riga.\newline
    Fine (riga) \\ \hline

    \textbf{Scenario alternativo} & 
    7b. Precondizione: un campo specifico inline non è valido.\newline
    Il sistema registra un warning per quel campo specifico e lo salta; la categoria viene creata senza di esso.\newline
    Fine (riga) \\ \hline

    \textbf{Scenario alternativo} & 
    7c. Precondizione: errore di salvataggio (I/O) della categoria.\newline
    Il sistema esegue il rollback in memoria, registra un errore e prosegue.\newline
    Fine (riga) \\ \hline

    \textbf{Scenario alternativo} & 
    8a. Precondizione: la categoria indicata nella riga PROPOSTA non esiste.\newline
    Il sistema registra un errore e salta la riga.\newline
    Fine (riga) \\ \hline

    \textbf{Scenario alternativo} & 
    8b. Precondizione: token campo = valore non validi o campo non esistente.\newline
    Il sistema registra un warning per ogni token non riconosciuto e prosegue con i restanti campi.\newline
    Fine (riga) \\ \hline

    \textbf{Scenario alternativo} & 
    8c. Precondizione: la proposta risulta non valida dopo la validazione.\newline
    Il sistema registra un errore con i problemi rilevati. La proposta non viene pubblicata.\newline
    Fine (riga) \\ \hline

    \textbf{Scenario alternativo} & 
    8d. Precondizione: errore di salvataggio (I/O) della proposta pubblicata.\newline
    Il sistema esegue il rollback in memoria, registra un errore e prosegue.\newline
    Fine (riga) \\ \hline
\end{tabularx}

\renewcommand{\arraystretch}{1.4} 

\begin{tabularx}{\textwidth}{|l|X|}
    \hline
    \textbf{Nome} & \textbf{UC30 – Importa più file batch in sequenza} \\ \hline
    \textbf{Attore} & Configuratore \\ \hline
    \endfirsthead 

    \hline
    \multicolumn{2}{|l|}{\textit{Continua dalla pagina precedente...}} \\ \hline
    \endhead 

    \hline
    \multicolumn{2}{|r|}{\textit{Continua nella pagina successiva...}} \\ \hline
    \endfoot 

    \hline
    \endlastfoot

    \textbf{Scenario principale} & 
    Precondizione: Il configuratore è autenticato.\newline
    Trigger: Selezione dell'opzione "Importa da file batch" $\rightarrow$ "Importa più file in sequenza".\newline
    1. Il sistema chiede di inserire i percorsi dei file batch uno alla volta.\newline
    2. Il configuratore inserisce i percorsi (uno per invio).\newline
    3. Il configuratore preme INVIO con il campo vuoto per confermare la lista.\newline
    4. Il sistema mostra il numero di file selezionati e avvia l'elaborazione sequenziale.\newline
    5. Per ciascun file della lista, il sistema esegue l'elaborazione descritta in UC29 (passi 3–9).\newline
    6. I risultati parziali di ogni file vengono aggregati in un unico \textit{BatchRisultato} complessivo tramite il metodo \textit{aggiungi(parziale)}.\newline
    7. Al termine di tutti i file, il sistema mostra il resoconto aggregato con il numero totale di righe elaborate, successi, warning ed errori, e l'elenco dei file processati.\newline
    Fine \\ \hline
    
    \textbf{Scenario alternativo} & 
    2a. Precondizione: l'utente digita "annulla" invece di un percorso.\newline
    Il sistema annulla l'intera operazione senza apportare modifiche.\newline
    Fine \\ \hline

    \textbf{Scenario alternativo} & 
    3a. Precondizione: l'utente preme INVIO senza aver inserito alcun percorso.\newline
    Il sistema mostra un errore "Nessun file specificato" e interrompe l'operazione.\newline
    Fine \\ \hline

    \textbf{Scenario alternativo} & 
    5a. Precondizione: uno dei file indicati non esiste o non è leggibile (IOException).\newline
    Il sistema registra un errore nel resoconto aggregato (indicando il file) e continua con il file successivo nella lista.\newline
    Fine (file) \\ \hline

    \textbf{Scenario alternativo} & 
    5b. Precondizione: errori su singole righe di un file (comandi errati, duplicati, validazione fallita).\newline
    Si applicano le stesse alternative di UC29 (5a–8d) con i relativi rollback per riga. L'elaborazione prosegue per le righe e i file successivi.\newline
    Fine \\ \hline
\end{tabularx}
\section{Casi d'uso grafici}

\begin{figure}[htbp]
    \centering
    \makebox[\textwidth][c]{
        \includegraphics[width=1.3\textwidth]{immagini/casi uso grafici v5 definitivo mini.png} 
    }
    \label{fig:usecase_v5}
\end{figure}

\clearpage
\section{Diagramma delle classi}
\clearpage
\begin{figure}[p]
    \centering
    \includegraphics[width=\textwidth]{immagini/uml v5.png}
    \label{fig:diagramma_classi_v5_view_controller_persistence}
\end{figure}

\begin{figure}[p]
    \centering
    \includegraphics[width=0.6\textwidth]{immagini/uml v5 2.png}
    \label{fig:diagramma_classi_v5_model}
\end{figure}

% --- CAPITOLO 6 ---
\chapter{Versione Finale}
\section{Casi d'uso testuali}

\begin{table}[htbp]
    \centering
    \renewcommand{\arraystretch}{1.4} 
    \begin{tabularx}{\textwidth}{|l|X|}
    \hline
    \textbf{Nome}               & \textbf{UC01 - Avvio e Login} \\ \hline
    \textbf{Attore}             & Configuratore \\ \hline
    \textbf{Scenario principale} & 
    Precondizione: L'applicazione è avviata. Transizioni automatiche (UC20) completate dal sistema.\newline
    Trigger: L'utente seleziona "Configuratore" dal menu iniziale.\newline
    1. Il sistema mostra la schermata di login dedicata.\newline
    2. Il configuratore inserisce username e password.\newline
    3. Il sistema verifica la corrispondenza delle credenziali.\newline
    4. Il sistema autentica il configuratore.\newline
    5. Il sistema verifica se si tratta del primissimo accesso assoluto (credenziali di default).\newline
    6. <<include>> "UC02 - Imposta credenziali personali" (se primo accesso).\newline
    7. Il sistema mostra il menu principale aggiornato coi contatori in tempo reale.\newline
    Postcondizione: Il configuratore è autenticato nel back-end.\newline
    Fine \\ \hline
    \textbf{Scenario alternativo} & 
    2a. Precondizione: il configuratore digita "esci" come username.\newline
    Il sistema interrompe l'operazione e torna al menu di selezione ruolo.\newline
    Fine \\ \hline
    \textbf{Scenario alternativo} & 
    3a. Precondizione: le credenziali non sono valide.\newline
    Il sistema mostra un errore e indica i tentativi rimasti.\newline
    Torna al punto 2 (se tentativi $< 3$), altrimenti il sistema blocca l'accesso alla sessione.\newline
    Fine \\ \hline
    \end{tabularx}
    \label{tab:uc01_v3}
\end{table}

\begin{tabularx}{\textwidth}{|l|X|}
    \hline
    \textbf{Nome}               & \textbf{UC02 - Imposta credenziali personali} \\ \hline
    \textbf{Attore}             & Configuratore \\ \hline
    \endfirsthead 
    \hline
    \multicolumn{2}{|l|}{\textit{Continua dalla pagina precedente...}} \\ \hline
    \endhead 
    \hline
    \multicolumn{2}{|r|}{\textit{Continua nella pagina successiva...}} \\ \hline
    \endfoot 
    \hline
    \endlastfoot 
    \textbf{Scenario principale} & 
    Precondizione: Login effettuato con credenziali di default (primo accesso).\newline
    Trigger: Il sistema rileva la necessità di cambiare le credenziali.\newline
    1. Il sistema richiede di scegliere le credenziali personali.\newline
    2. Il configuratore inserisce il nuovo username.\newline
    3. Il sistema verifica che lo username non coincida con la parola riservata "esci".\newline
    4. Il configuratore inserisce la nuova password.\newline
    5. Il configuratore inserisce la password per conferma.\newline
    6. Il sistema verifica che le due password coincidano.\newline
    7. Il sistema verifica che username e password non siano stringhe vuote/blank.\newline
    8. Il sistema verifica l'unicità globale dello username (tra Configuratori e Fruitori).\newline
    9. Il sistema aggiorna le credenziali in memoria.\newline
    10. Il sistema salva i dati in modo persistente.\newline
    11. Il sistema mostra la conferma e accede al menu.\newline
    Postcondizione: Le credenziali di default sono sostituite e salvate in modo permanente.\newline
    Fine \\ \hline
    \textbf{Scenario alternativo} & 
    2a. Precondizione: il configuratore digita "annulla" nel campo username.\newline
    Il sistema annulla l'operazione ed esegue il logout automatico.\newline
    Fine \\ \hline
    \textbf{Scenario alternativo} & 
    3a. Precondizione: lo username inserito è "esci".\newline
    Il sistema mostra un errore, annulla l'operazione ed esegue il logout.\newline
    Fine \\ \hline
    \textbf{Scenario alternativo} & 
    4a. Precondizione: il configuratore digita "annulla" nel campo password.\newline
    Il sistema annulla l'operazione ed esegue il logout automatico.\newline
    Fine \\ \hline
    \textbf{Scenario alternativo} & 
    6a. Precondizione: le password non coincidono.\newline
    Il sistema mostra un errore, annulla l'operazione ed esegue il logout.\newline
    Fine \\ \hline
    \textbf{Scenario alternativo} & 
    7a. Precondizione: campi lasciati vuoti.\newline
    Il sistema mostra un errore, annulla l'operazione ed esegue il logout.\newline
    Fine \\ \hline
    \textbf{Scenario alternativo} & 
    8a. Precondizione: username già in uso nel sistema.\newline
    Il sistema mostra un errore, annulla l'operazione ed esegue il logout.\newline
    Fine \\ \hline
    \textbf{Scenario alternativo} & 
    10a. Precondizione: Errore di salvataggio (I/O).\newline
    Il sistema esegue un rollback in memoria, mostra un errore ed esegue il logout.\newline
    Fine \\ \hline
\end{tabularx}

\begin{tabularx}{\textwidth}{|l|X|}
    \hline
    \textbf{Nome}               & \textbf{UC03 - Logout} \\ \hline
    \textbf{Attore}             & Configuratore \\ \hline
    \endfirsthead 

    \hline
    \multicolumn{2}{|l|}{\textit{Continua dalla pagina precedente...}} \\ \hline
    \endhead 

    \hline
    \multicolumn{2}{|r|}{\textit{Continua nella pagina successiva...}} \\ \hline
    \endfoot 

    \hline
    \endlastfoot

    \textbf{Scenario principale} & 
    Precondizione: Configuratore autenticato. \newline
    Trigger: Selezione "Logout". \newline
    1. Il configuratore seleziona Logout. \newline
    2. Il sistema scarta le proposte non pubblicate della sessione corrente (stato BOZZA o VALIDA).\newline
    3. Se erano presenti proposte non pubblicate, il sistema avvisa il configuratore.\newline
    4. Il sistema termina la sessione corrente. \newline
    5. Il sistema chiede se accedere con un altro account. \newline
    6. Sì: ritorna a UC01. \newline
    7. No: termina l'applicazione. \newline
    Fine \\ \hline
    
    \textbf{Scenario alternativo} & 
    5a. Precondizione: Il configuratore inserisce una risposta non valida (diversa da "s" o "n").\newline
    Il sistema mostra un messaggio di errore testuale e ripropone la domanda. \\ \hline
\end{tabularx}

\begin{tabularx}{\textwidth}{|l|X|}
    \hline
    \textbf{Nome}               & \textbf{UC04 - Visualizza campi base} \\ \hline
    \textbf{Attore}             & Configuratore \\ \hline
    \endfirsthead 

    \hline
    \multicolumn{2}{|l|}{\textit{Continua dalla pagina precedente...}} \\ \hline
    \endhead 

    \hline
    \multicolumn{2}{|r|}{\textit{Continua nella pagina successiva...}} \\ \hline
    \endfoot 

    \hline
    \endlastfoot

    \textbf{Scenario principale} & 
    Precondizione: Il configuratore è autenticato. I campi base sono stati inizializzati.\newline
    Trigger: Il configuratore seleziona "Gestione campi base" dal menu principale.\newline
    1. Il configuratore seleziona la voce "Gestione campi base".\newline
    2. Il sistema mostra l'elenco dei campi base con nome e tipo, evidenziando che sono tutti obbligatori (Titolo, Numero partecipanti, ecc.).\newline
    3. Il sistema indica che i campi base sono immutabili.\newline
    4. Il configuratore preme INVIO per tornare al menu principale.\newline
    Fine \\ \hline
    
    \textbf{Scenario alternativo} & 
    2a. Precondizione: Campi base non inizializzati.\newline
    Il sistema li inizializza automaticamente al primo avvio e poi mostra l'elenco.\newline
    Fine \\ \hline
\end{tabularx}

\renewcommand{\arraystretch}{1.4} 

\begin{tabularx}{\textwidth}{|l|X|}
    \hline
    \textbf{Nome}               & \textbf{UC05 - Aggiungi campo comune} \\ \hline
    \textbf{Attore}             & Configuratore \\ \hline
    \endfirsthead 

    \hline
    \multicolumn{2}{|l|}{\textit{Continua dalla pagina precedente...}} \\ \hline
    \endhead 

    \hline
    \multicolumn{2}{|r|}{\textit{Continua nella pagina successiva...}} \\ \hline
    \endfoot 

    \hline
    \endlastfoot

    \textbf{Scenario principale} & 
    Precondizione: Il configuratore è autenticato. \newline
    Trigger: Il configuratore seleziona "Aggiungi campo comune" nel menu "Gestione campi comuni". \newline
    1. Il configuratore seleziona "Aggiungi campo comune". \newline
    2. Il sistema chiede il nome del nuovo campo. \newline
    3. Il configuratore inserisce il nome. \newline
    4. Il sistema chiede se il campo è obbligatorio o facoltativo. \newline
    5. Il configuratore risponde (s/n). \newline
    6. Il sistema verifica che il nome non coincida con un campo base o comune esistente. \newline
    7. Il sistema aggiunge il campo comune e salva in modo persistente. \newline
    8. Il sistema mostra conferma dell'operazione. \newline
    Fine \\ \hline
    
    \textbf{Scenario alternativo} & 
    3a. Precondizione: Nome vuoto. \newline
    Il sistema mostra errore e richiede nuovamente il nome. \\ \hline

    \textbf{Scenario alternativo} & 
    6a. Precondizione: Conflitto con campo base. \newline
    Il sistema mostra errore e non aggiunge il campo. \newline
    Fine \\ \hline

    \textbf{Scenario alternativo} & 
    6b. Precondizione: Conflitto con campo comune esistente. \newline
    Il sistema mostra errore e non aggiunge il campo. \newline
    Fine \\ \hline

    \textbf{Scenario alternativo} & 
    7a. Precondizione: Errore di salvataggio (I/O). \newline
    Il sistema aggiunge il campo in memoria ma avvisa dell'errore di salvataggio fisico. \newline
    Fine \\ \hline
\end{tabularx}

\renewcommand{\arraystretch}{1.4} 

\begin{tabularx}{\textwidth}{|l|X|}
    \hline
    \textbf{Nome}               & \textbf{UC06 - Rimuovi campo comune} \\ \hline
    \textbf{Attore}             & Configuratore \\ \hline
    \endfirsthead 

    \hline
    \multicolumn{2}{|l|}{\textit{Continua dalla pagina precedente...}} \\ \hline
    \endhead 

    \hline
    \multicolumn{2}{|r|}{\textit{Continua nella pagina successiva...}} \\ \hline
    \endfoot 

    \hline
    \endlastfoot

    \textbf{Scenario principale} & 
    Precondizione: Il configuratore è autenticato. Esiste almeno un campo comune. \newline
    Trigger: Il configuratore seleziona "Rimuovi campo comune". \newline
    1. Il configuratore seleziona "Rimuovi campo comune". \newline
    2. Il sistema mostra l'elenco dei campi comuni esistenti. \newline
    3. Il configuratore inserisce il nome del campo da rimuovere. \newline
    4. Il sistema rimuove il campo comune e salva in modo persistente. \newline
    5. Il sistema mostra conferma dell'operazione. \newline
    Fine \\ \hline
    
    \textbf{Scenario alternativo} & 
    3a. Precondizione: Campo non trovato. \newline
    Il sistema mostra errore e non esegue la rimozione. \newline
    Fine \\ \hline

    \textbf{Scenario alternativo} & 
    4a. Precondizione: Errore di salvataggio (I/O). \newline
    Il sistema rimuove il campo in memoria ma avvisa dell'errore di salvataggio fisico. \newline
    Fine \\ \hline
\end{tabularx}

\renewcommand{\arraystretch}{1.4} 

\begin{tabularx}{\textwidth}{|l|X|}
    \hline
    \textbf{Nome}               & \textbf{UC07 - Modifica obbligatorietà campo comune} \\ \hline
    \textbf{Attore}             & Configuratore \\ \hline
    \endfirsthead 

    \hline
    \multicolumn{2}{|l|}{\textit{Continua dalla pagina precedente...}} \\ \hline
    \endhead 

    \hline
    \multicolumn{2}{|r|}{\textit{Continua nella pagina successiva...}} \\ \hline
    \endfoot 

    \hline
    \endlastfoot

    \textbf{Scenario principale} & 
    Precondizione: Il configuratore è autenticato. Esiste almeno un campo comune. \newline
    Trigger: Il configuratore seleziona "Modifica obbligatorietà campo comune". \newline
    1. Il configuratore seleziona "Modifica obbligatorietà campo comune". \newline
    2. Il sistema mostra l'elenco dei campi comuni con l'obbligatorietà attuale. \newline
    3. Il configuratore inserisce il nome del campo da modificare. \newline
    4. Il configuratore indica il nuovo valore di obbligatorietà (s/n). \newline
    5. Il sistema aggiorna il campo e salva in modo persistente. \newline
    6. Il sistema mostra conferma dell'operazione. \newline
    Fine \\ \hline
    
    \textbf{Scenario alternativo} & 
    3a. Precondizione: Campo non trovato. \newline
    Il sistema mostra errore e non applica modifiche. \newline
    Fine \\ \hline

    \textbf{Scenario alternativo} & 
    5a. Precondizione: Errore di salvataggio (I/O). \newline
    Il sistema aggiorna il campo in memoria ma avvisa dell'errore di salvataggio fisico. \newline
    Fine \\ \hline
\end{tabularx}

\renewcommand{\arraystretch}{1.4} 

\begin{tabularx}{\textwidth}{|l|X|}
    \hline
    \textbf{Nome}               & \textbf{UC08 - Aggiungi categoria} \\ \hline
    \textbf{Attore}             & Configuratore \\ \hline
    \endfirsthead 

    \hline
    \multicolumn{2}{|l|}{\textit{Continua dalla pagina precedente...}} \\ \hline
    \endhead 

    \hline
    \multicolumn{2}{|r|}{\textit{Continua nella pagina successiva...}} \\ \hline
    \endfoot 

    \hline
    \endlastfoot

    \textbf{Scenario principale} & 
    Precondizione: Il configuratore è autenticato. I campi base sono inizializzati. \newline
    Trigger: Il configuratore seleziona "Aggiungi categoria". \newline
    1. Il configuratore seleziona "Aggiungi categoria". \newline
    2. Il sistema chiede il nome della nuova categoria. \newline
    3. Il configuratore inserisce il nome. \newline
    4. Il sistema verifica che il nome non sia già in uso. \newline
    5. Il sistema crea la categoria e salva in modo persistente. \newline
    6. Il sistema chiede se aggiungere subito campi specifici. \newline
    7. Sì: il caso d'uso estende UC10 (aggiunta campi specifici). \newline
    8. No: il sistema torna al menu categorie. \newline
    9. Il sistema mostra conferma dell'operazione. \newline
    Fine \\ \hline
    
    \textbf{Scenario alternativo} & 
    3a. Precondizione: Nome vuoto. \newline
    Errore e richiesta nuovo inserimento. \\ \hline

    \textbf{Scenario alternativo} & 
    4a. Precondizione: Categoria duplicata. \newline
    Errore e nessuna creazione. \newline
    Fine \\ \hline

    \textbf{Scenario alternativo} & 
    5a. Precondizione: Errore di salvataggio (I/O). \newline
    Il sistema crea la categoria in memoria ma avvisa dell'errore di salvataggio fisico. \newline
    Fine \\ \hline
\end{tabularx}

\renewcommand{\arraystretch}{1.4} 

\begin{tabularx}{\textwidth}{|l|X|}
    \hline
    \textbf{Nome}               & \textbf{UC09 - Rimuovi categoria} \\ \hline
    \textbf{Attore}             & Configuratore \\ \hline
    \endfirsthead 

    \hline
    \multicolumn{2}{|l|}{\textit{Continua dalla pagina precedente...}} \\ \hline
    \endhead 

    \hline
    \multicolumn{2}{|r|}{\textit{Continua nella pagina successiva...}} \\ \hline
    \endfoot 

    \hline
    \endlastfoot

    \textbf{Scenario principale} & 
    Precondizione: Il configuratore è autenticato. Esiste almeno una categoria. \newline
    Trigger: Il configuratore seleziona "Rimuovi categoria". \newline
    1. Il configuratore seleziona "Rimuovi categoria". \newline
    2. Il sistema mostra l'elenco categorie. \newline
    3. Il configuratore inserisce il nome della categoria da rimuovere. \newline
    4. Il sistema chiede conferma esplicita prima della rimozione. \newline
    5. Il configuratore conferma. \newline
    6. Il sistema rimuove la categoria e tutti i suoi campi specifici. \newline
    7. Il sistema salva i dati in modo persistente. \newline
    8. Il sistema mostra conferma. \newline
    Fine \\ \hline
    
    \textbf{Scenario alternativo} & 
    3a. Precondizione: Categoria non trovata. \newline
    Errore, nessuna rimozione. \newline
    Fine \\ \hline

    \textbf{Scenario alternativo} & 
    5a. Precondizione: Annulla. \newline
    Il configuratore non conferma; operazione interrotta. \newline
    Fine \\ \hline

    \textbf{Scenario alternativo} & 
    7a. Precondizione: Errore di salvataggio (I/O). \newline
    Il sistema rimuove la categoria in memoria ma avvisa dell'errore di salvataggio fisico. \newline
    Fine \\ \hline
\end{tabularx}

\renewcommand{\arraystretch}{1.4} 

\begin{tabularx}{\textwidth}{|l|X|}
    \hline
    \textbf{Nome}               & \textbf{UC10 - Aggiungi campo specifico a categoria} \\ \hline
    \textbf{Attore}             & Configuratore \\ \hline
    \endfirsthead 

    \hline
    \multicolumn{2}{|l|}{\textit{Continua dalla pagina precedente...}} \\ \hline
    \endhead 

    \hline
    \multicolumn{2}{|r|}{\textit{Continua nella pagina successiva...}} \\ \hline
    \endfoot 

    \hline
    \endlastfoot

    \textbf{Scenario principale} & 
    Precondizione: Il configuratore è autenticato. Esiste almeno una categoria. \newline
    Trigger: Il configuratore sceglie di aggiungere un campo specifico a una categoria. \newline
    1. Il configuratore seleziona "Gestisci campi specifici" e indica la categoria. \newline
    2. Il sistema verifica che la categoria esista. \newline
    3. Il configuratore seleziona "Aggiungi campo specifico". \newline
    4. Il sistema chiede il nome del nuovo campo. \newline
    5. Il configuratore inserisce il nome. \newline
    6. Il sistema chiede se il campo è obbligatorio o facoltativo. \newline
    7. Il configuratore risponde (s/n). \newline
    8. Il sistema verifica che il nome non coincida con un campo base, comune o specifico già presente. \newline
    9. Il sistema aggiunge il campo specifico alla categoria e salva. \newline
    10. Il sistema mostra conferma e permette inserimenti successivi (INVIO per terminare).\newline
    Fine \\ \hline
    
    \textbf{Scenario alternativo} & 
    2a. Precondizione: Categoria non esiste. \newline
    Il sistema mostra un messaggio di errore. \newline
    Fine \\ \hline

    \textbf{Scenario alternativo} & 
    5a. Precondizione: Nome vuoto. \newline
    Il sistema mostra un errore e richiede un nuovo inserimento. \\ \hline

    \textbf{Scenario alternativo} & 
    8a. Precondizione: Nome duplicato. \newline
    Il sistema mostra un errore e non procede con l'aggiunta. \newline
    Fine \\ \hline

    \textbf{Scenario alternativo} & 
    9a. Precondizione: Errore di salvataggio (I/O). \newline
    Il sistema aggiunge il campo in memoria ma avvisa dell'errore di salvataggio fisico. \newline
    Fine \\ \hline
\end{tabularx}

\renewcommand{\arraystretch}{1.4} 

\begin{tabularx}{\textwidth}{|l|X|}
    \hline
    \textbf{Nome}               & \textbf{UC11 - Rimuovi campo specifico da categoria} \\ \hline
    \textbf{Attore}             & Configuratore \\ \hline
    \endfirsthead 

    \hline
    \multicolumn{2}{|l|}{\textit{Continua dalla pagina precedente...}} \\ \hline
    \endhead 

    \hline
    \multicolumn{2}{|r|}{\textit{Continua nella pagina successiva...}} \\ \hline
    \endfoot 

    \hline
    \endlastfoot

    \textbf{Scenario principale} & 
    Precondizione: Il configuratore è autenticato. Esiste una categoria con almeno un campo specifico. \newline
    Trigger: Il configuratore seleziona "Rimuovi campo specifico". \newline
    1. Il configuratore seleziona la categoria. \newline
    2. Il sistema mostra i campi specifici della categoria. \newline
    3. Il configuratore inserisce il nome del campo da rimuovere. \newline
    4. Il sistema rimuove il campo specifico e salva in modo persistente. \newline
    5. Il sistema mostra conferma. \newline
    Fine \\ \hline
    
    \textbf{Scenario alternativo} & 
    3a. Precondizione: Campo non trovato. \newline
    Il sistema mostra un messaggio di errore. \newline
    Fine \\ \hline

    \textbf{Scenario alternativo} & 
    4a. Precondizione: Errore di salvataggio (I/O). \newline
    Il sistema rimuove il campo in memoria ma avvisa dell'errore di salvataggio fisico. \newline
    Fine \\ \hline
\end{tabularx}

\renewcommand{\arraystretch}{1.4} 

\begin{tabularx}{\textwidth}{|l|X|}
    \hline
    \textbf{Nome}               & \textbf{UC12 - Modifica obbligatorietà campo specifico} \\ \hline
    \textbf{Attore}             & Configuratore \\ \hline
    \endfirsthead 

    \hline
    \multicolumn{2}{|l|}{\textit{Continua dalla pagina precedente...}} \\ \hline
    \endhead 

    \hline
    \multicolumn{2}{|r|}{\textit{Continua nella pagina successiva...}} \\ \hline
    \endfoot 

    \hline
    \endlastfoot

    \textbf{Scenario principale} & 
    Precondizione: Il configuratore è autenticato. Esiste una categoria con almeno un campo specifico. \newline
    Trigger: Il configuratore seleziona "Modifica obbligatorietà campo specifico". \newline
    1. Il configuratore seleziona la categoria. \newline
    2. Il sistema mostra i campi specifici con obbligatorietà attuale. \newline
    3. Il configuratore inserisce il nome del campo da modificare. \newline
    4. Il configuratore indica il nuovo valore (s/n). \newline
    5. Il sistema aggiorna il campo specifico e salva in modo persistente. \newline
    6. Il sistema mostra conferma. \newline
    Fine \\ \hline
    
    \textbf{Scenario alternativo} & 
    3a. Precondizione: Campo non trovato. \newline
    Il sistema mostra un messaggio di errore. \newline
    Fine \\ \hline

    \textbf{Scenario alternativo} & 
    5a. Precondizione: Errore di salvataggio (I/O). \newline
    Il sistema aggiorna il campo in memoria ma avvisa dell'errore di salvataggio fisico. \newline
    Fine \\ \hline
\end{tabularx}

\renewcommand{\arraystretch}{1.4} 

\begin{tabularx}{\textwidth}{|l|X|}
    \hline
    \textbf{Nome}               & \textbf{UC13 - Visualizza categorie e campi} \\ \hline
    \textbf{Attore}             & Configuratore \\ \hline
    \endfirsthead 

    \hline
    \multicolumn{2}{|l|}{\textit{Continua dalla pagina precedente...}} \\ \hline
    \endhead 

    \hline
    \multicolumn{2}{|r|}{\textit{Continua nella pagina successiva...}} \\ \hline
    \endfoot 

    \hline
    \endlastfoot

    \textbf{Scenario principale} & 
    Precondizione: Il configuratore è autenticato. \newline
    Trigger: Il configuratore seleziona "Visualizza riepilogo categorie e campi". \newline
    1. Il configuratore seleziona "Visualizza riepilogo". \newline
    2. Il sistema mostra i campi BASE. \newline
    3. Il sistema mostra i campi COMUNI (con indicazione di obbligatorio/facoltativo). \newline
    4. Per ogni categoria, il sistema mostra: nome categoria e campi specifici (con indicazione di obbligatorio/facoltativo). \newline
    5. Il configuratore preme INVIO per tornare al menu. \newline
    Fine \\ \hline
    
    \textbf{Scenario alternativo} & 
    3a. Precondizione: Nessun campo comune definito. \newline
    Il sistema mostra la dicitura "(nessuno definito)". \\ \hline

    \textbf{Scenario alternativo} & 
    4a. Precondizione: Nessuna categoria definita nel sistema. \newline
    Il sistema mostra la dicitura "(nessuna categoria definita)". \\ \hline

    \textbf{Scenario alternativo} & 
    4b. Precondizione: Una categoria non possiede campi specifici. \newline
    Il sistema mostra la dicitura "(nessun campo specifico)" sotto il nome della categoria interessata. \\ \hline
\end{tabularx}

\renewcommand{\arraystretch}{1.4} 

\begin{tabularx}{\textwidth}{|l|X|}
    \hline
    \textbf{Nome}               & \textbf{UC14 - Visualizza dettaglio categoria} \\ \hline
    \textbf{Attore}             & Configuratore \\ \hline
    \endfirsthead 

    \hline
    \multicolumn{2}{|l|}{\textit{Continua dalla pagina precedente...}} \\ \hline
    \endhead 

    \hline
    \multicolumn{2}{|r|}{\textit{Continua nella pagina successiva...}} \\ \hline
    \endfoot 

    \hline
    \endlastfoot

    \textbf{Scenario principale} & 
    Precondizione: Il configuratore è autenticato. Esiste almeno una categoria. \newline
    Trigger: Il configuratore seleziona "Visualizza dettaglio categoria". \newline
    1. Il configuratore seleziona "Visualizza dettaglio categoria". \newline
    2. Il sistema richiede l'inserimento del nome della categoria. \newline
    3. Il configuratore inserisce il nome richiesto. \newline
    4. Il sistema mostra, per la categoria selezionata, il set completo dei dati: tutti i campi BASE, i campi COMUNI e i campi SPECIFICI della categoria stessa. \newline
    5. Il configuratore preme INVIO per terminare la visualizzazione e tornare al menu. \newline
    Fine \\ \hline
    
    \textbf{Scenario alternativo} & 
    3a. Precondizione: La categoria inserita non è presente nell'archivio. \newline
    Il sistema mostra un messaggio di errore e l'operazione si interrompe. \newline
    Fine \\ \hline
\end{tabularx}

\begin{tabularx}{\textwidth}{|l|X|}
    \hline
    \textbf{Nome}               & \textbf{UC15 - Crea proposta} \\ \hline
    \textbf{Attore}             & Configuratore \\ \hline
    \endfirsthead 

    \hline
    \multicolumn{2}{|l|}{\textit{Continua dalla pagina precedente...}} \\ \hline
    \endhead 

    \hline
    \multicolumn{2}{|r|}{\textit{Continua nella pagina successiva...}} \\ \hline
    \endfoot 

    \hline
    \endlastfoot

    \textbf{Scenario principale} & 
    Precondizione: Il configuratore è autenticato ed esiste almeno una categoria.\newline
    Trigger: Selezione "Nuova proposta" dal menu gestionale.\newline
    1. Il sistema mostra l'elenco delle categorie disponibili.\newline
    2. Il configuratore inserisce il nome della categoria desiderata.\newline
    3. Il sistema verifica l'esistenza della categoria.\newline
    4. Il sistema cattura lo snapshot dei campi (base, comuni e specifici) rendendoli indipendenti da future modifiche alla configurazione.\newline
    5. Il sistema assegna un ID univoco progressivo estratto dal contatore globale.\newline
    6. Il sistema aggiunge la proposta alla sessione corrente in stato BOZZA.\newline
    7. <<include>> "UC16 - Compila campi proposta".\newline
    8. Il sistema mostra il riepilogo e, se in stato VALIDA, offre la pubblicazione immediata (estende UC17).\newline
    Fine \\ \hline
    
    \textbf{Scenario alternativo} & 
    1a. Precondizione: l'archivio categorie è vuoto.\newline
    Il sistema mostra l'errore "Nessuna categoria disponibile" e impedisce la creazione.\newline
    Fine \\ \hline

    \textbf{Scenario alternativo} & 
    3a. Precondizione: categoria non trovata in base all'input dell'utente.\newline
    Il sistema mostra un errore e torna al menu proposte.\newline
    Fine \\ \hline
\end{tabularx}

\begin{tabularx}{\textwidth}{|l|X|}
    \hline
    \textbf{Nome}               & \textbf{UC16 - Compila / Modifica campi proposta} \\ \hline
    \textbf{Attore}             & Configuratore \\ \hline
    \endfirsthead 

    \hline
    \multicolumn{2}{|l|}{\textit{Continua dalla pagina precedente...}} \\ \hline
    \endhead 

    \hline
    \multicolumn{2}{|r|}{\textit{Continua nella pagina successiva...}} \\ \hline
    \endfoot 

    \hline
    \endlastfoot

    \textbf{Scenario principale} & 
    Precondizione: Esiste almeno una proposta in sessione.\newline
    Trigger: Estensione da UC15 o selezione "Continua compilazione".\newline
    1. Il sistema presenta i campi dello snapshot uno alla volta.\newline
    2. Per ogni campo, il sistema mostra l'etichetta [OBB/FAC], il formato richiesto e il valore attuale.\newline
    3. Il configuratore inserisce un valore o preme INVIO per mantenere il valore attuale.\newline
    4. Il sistema valida l'input e ricalcola automaticamente lo stato (BOZZA/VALIDA).\newline
    5. Il sistema verifica i vincoli temporali: Termine iscrizione $>$ oggi; Data inizio $\ge$ Termine + 2gg; Data fine $\ge$ Data inizio; Formato Ora (HH:MM).\newline
    6. Il sistema mostra il riepilogo finale e l'elenco degli errori se lo stato è BOZZA.\newline
    Fine \\ \hline
    
    \textbf{Scenario alternativo} & 
    4a. Precondizione: formato data (gg/mm/aaaa) o ora (HH:MM) errato.\newline
    Il sistema mostra un errore specifico e non aggiorna il valore del campo. Prosegue con il campo successivo. \\ \hline
\end{tabularx}

\begin{tabularx}{\textwidth}{|l|X|}
    \hline
    \textbf{Nome}               & \textbf{UC17 - Pubblica proposta in bacheca} \\ \hline
    \textbf{Attore}             & Configuratore \\ \hline
    \endfirsthead 
    \hline
    \multicolumn{2}{|l|}{\textit{Continua dalla pagina precedente...}} \\ \hline
    \endhead 
    \hline
    \multicolumn{2}{|r|}{\textit{Continua nella pagina successiva...}} \\ \hline
    \endfoot 
    \hline
    \endlastfoot
    \textbf{Scenario principale} & 
    Precondizione: Il configuratore è autenticato e c'è almeno una proposta VALIDA in sessione.\newline
    Trigger: Selezione "Pubblica proposta in bacheca".\newline
    1. Il sistema mostra l'elenco delle proposte in sessione attualmente VALIDE.\newline
    2. Il configuratore inserisce l'ID della proposta da pubblicare.\newline
    3. Il sistema cattura la data odierna.\newline
    4. Il sistema ricalcola la validità della proposta in base alla data odierna.\newline
    5. Il sistema imposta lo stato su APERTA e registra la data nello storico.\newline
    6. Il sistema sposta la proposta dalla sessione all'Archivio globale.\newline
    7. Il sistema salva le modifiche in modo persistente.\newline
    8. Il sistema mostra conferma dell'operazione.\newline
    Postcondizione: La proposta ha stato APERTA, è salvata nell'archivio ed è visibile ai fruitori.\newline
    Fine \\ \hline
    \textbf{Scenario alternativo} & 
    1a. Precondizione: nessuna proposta in sessione è in stato VALIDA.\newline
    Il sistema mostra un errore.\newline
    Fine \\ \hline
    \textbf{Scenario alternativo} & 
    2a. Precondizione: l'input non è numerico o non corrisponde a una proposta in elenco.\newline
    Il sistema mostra un errore.\newline
    Fine \\ \hline
    \textbf{Scenario alternativo} & 
    4a. Precondizione: la validità è scaduta a causa del cambio di data (es. superata mezzanotte).\newline
    Il sistema nega la pubblicazione e mostra l'elenco degli errori.\newline
    Fine \\ \hline
    \textbf{Scenario alternativo} & 
    7a. Precondizione: Errore di salvataggio (I/O).\newline
    Il sistema esegue un rollback completo riportando la proposta allo stato VALIDA in sessione.\newline
    Fine \\ \hline
\end{tabularx}

\begin{tabularx}{\textwidth}{|l|X|}
    \hline
    \textbf{Nome}               & \textbf{UC18 - Elimina proposta dalla sessione} \\ \hline
    \textbf{Attore}             & Configuratore \\ \hline
    \endfirsthead 

    \hline
    \multicolumn{2}{|l|}{\textit{Continua dalla pagina precedente...}} \\ \hline
    \endhead 

    \hline
    \multicolumn{2}{|r|}{\textit{Continua nella pagina successiva...}} \\ \hline
    \endfoot 

    \hline
    \endlastfoot

    \textbf{Scenario principale} & 
    Precondizione: Il configuratore è autenticato. Esiste almeno una proposta in sessione (stato BOZZA o VALIDA).\newline
    Trigger: Il configuratore seleziona "Elimina proposta" nel menu proposte.\newline
    1. Il sistema mostra l'elenco delle proposte in sessione.\newline
    2. Il configuratore inserisce l'ID della proposta da eliminare.\newline
    3. Il sistema chiede conferma esplicita prima dell'eliminazione.\newline
    4. Il configuratore conferma.\newline
    5. Il sistema rimuove definitivamente la proposta dalla sessione.\newline
    6. Il sistema mostra conferma dell'avvenuta operazione.\newline
    Fine \\ \hline
    
    \textbf{Scenario alternativo} & 
    2a. Precondizione: l'ID inserito non è valido o non corrisponde a una proposta presente in sessione.\newline
    Il sistema mostra un messaggio di errore e interrompe l'operazione.\newline
    Fine \\ \hline

    \textbf{Scenario alternativo} & 
    4a. Precondizione: il configuratore annulla l'operazione (nega la conferma).\newline
    Il sistema interrompe l'operazione senza apportare modifiche e torna al menu.\newline
    Fine \\ \hline
\end{tabularx}

\begin{tabularx}{\textwidth}{|l|X|}
    \hline
    \textbf{Nome}               & \textbf{UC19 - Visualizza bacheca Configuratore} \\ \hline
    \textbf{Attore}             & Configuratore \\ \hline
    \endfirsthead 
    \hline
    \multicolumn{2}{|l|}{\textit{Continua dalla pagina precedente...}} \\ \hline
    \endhead 
    \hline
    \multicolumn{2}{|r|}{\textit{Continua nella pagina successiva...}} \\ \hline
    \endfoot 
    \hline
    \endlastfoot
    \textbf{Scenario principale} & 
    Precondizione: Il configuratore è autenticato.\newline
    Trigger: Selezione "Visualizza bacheca".\newline
    1. Il configuratore richiede la visualizzazione della bacheca.\newline
    2. Il sistema recupera tutte le proposte con stato APERTA.\newline
    3. Il sistema raggruppa dinamicamente le proposte per categoria (inclusa la voce "(categoria rimossa)" per le eventuali proposte orfane).\newline
    4. Il sistema mostra per ogni proposta: ID, dettagli, valori dei campi e il numero di iscritti attuali rispetto al massimo consentito.\newline
    5. Il configuratore preme INVIO per terminare.\newline
    Postcondizione: Nessuna modifica allo stato del sistema.\newline
    Fine \\ \hline
    \textbf{Scenario alternativo} & 
    2a. Precondizione: non ci sono proposte APERTE in archivio.\newline
    Il sistema mostra il messaggio "(Nessuna proposta aperta in bacheca)".\newline
    Fine \\ \hline
    \textbf{Scenario alternativo} & 
    3a. Precondizione: categoria di una proposta rimossa dopo la pubblicazione.\newline
    Il sistema visualizza la proposta sotto l'etichetta "(categoria rimossa)".\newline
    Fine \\ \hline
\end{tabularx}

\begin{tabularx}{\textwidth}{|l|X|}
    \hline
    \textbf{Nome}               & \textbf{UC20 - Esecuzione transizioni automatiche di stato} \\ \hline
    \textbf{Attore}             & Sistema \\ \hline
    \endfirsthead 
    \hline
    \multicolumn{2}{|l|}{\textit{Continua dalla pagina precedente...}} \\ \hline
    \endhead 
    \hline
    \multicolumn{2}{|r|}{\textit{Continua nella pagina successiva...}} \\ \hline
    \endfoot 
    \hline
    \endlastfoot
    \textbf{Scenario principale} & 
    Precondizione: L'applicazione viene avviata. L'archivio contiene proposte.\newline
    Trigger: Inizializzazione del sistema.\newline
    1. Il sistema analizza l'archivio delle proposte.\newline
    2. Per ogni proposta APERTA con "Termine ultimo" superato, il sistema verifica se iscritti $\ge$ minimo richiesto.\newline
    3. Il sistema transita la proposta in stato CONFERMATA.\newline
    4. Il sistema genera una notifica di conferma per ogni aderente iscritto.\newline
    5. Per ogni proposta CONFERMATA con "Data conclusiva" superata, il sistema la transita in stato CONCLUSA.\newline
    6. Il sistema salva i nuovi stati e le notifiche in modo persistente.\newline
    7. Il sistema mostra a schermo il numero totale di transizioni effettuate.\newline
    Postcondizione: Stati aggiornati alla data odierna e notifiche distribuite.\newline
    Fine \\ \hline
    \textbf{Scenario alternativo} & 
    2a. Precondizione: la proposta APERTA è scaduta ma gli iscritti sono insufficienti (o il campo "Numero massimo" è corrotto/negativo).\newline
    Il sistema transita la proposta in ANNULLATA e genera una notifica di annullamento per ogni aderente.\newline
    Torna al punto 5 \\ \hline
    \textbf{Scenario alternativo} & 
    6a. Precondizione: Errore di salvataggio (I/O).\newline
    Il sistema esegue un rollback ricaricando i dati precedenti dal disco per evitare la generazione di notifiche duplicate ai riavvii successivi.\newline
    Fine \\ \hline
\end{tabularx}

\renewcommand{\arraystretch}{1.4} 

\begin{tabularx}{\textwidth}{|l|X|}
    \hline
    \textbf{Nome}               & \textbf{UC21 - Registrazione Fruitore} \\ \hline
    \textbf{Attore}             & Fruitore \\ \hline
    \endfirsthead 
    \hline
    \multicolumn{2}{|l|}{\textit{Continua dalla pagina precedente...}} \\ \hline
    \endhead 
    \hline
    \multicolumn{2}{|r|}{\textit{Continua nella pagina successiva...}} \\ \hline
    \endfoot 
    \hline
    \endlastfoot
    \textbf{Scenario principale} & 
    Precondizione: L'utente non è autenticato.\newline
    Trigger: Selezione "Registrazione (nuovo account)" dal menu di accesso.\newline
    1. Il sistema mostra la schermata di registrazione.\newline
    2. Il fruitore inserisce lo username desiderato.\newline
    3. Il fruitore inserisce la password.\newline
    4. Il fruitore reinserisce la password per conferma.\newline
    5. Il sistema verifica che le due password coincidano.\newline
    6. Il sistema verifica che username e password non siano vuoti.\newline
    7. Il sistema verifica l'unicità globale dello username.\newline
    8. Il sistema crea il profilo Fruitore in memoria.\newline
    9. Il sistema salva i dati in modo persistente.\newline
    10. Il sistema autentica automaticamente il fruitore e mostra il benvenuto.\newline
    Postcondizione: Il fruitore è registrato, autenticato e con uno spazio notifiche vuoto.\newline
    Fine \\ \hline
    \textbf{Scenario alternativo} & 
    2a. Precondizione: il fruitore digita "annulla" come username.\newline
    Il sistema interrompe l'operazione e torna al menu.\newline
    Fine \\ \hline
    \textbf{Scenario alternativo} & 
    3a. Precondizione: il fruitore digita "annulla" come password.\newline
    Il sistema interrompe l'operazione e torna al menu.\newline
    Fine \\ \hline
    \textbf{Scenario alternativo} & 
    5a. Precondizione: le password non coincidono.\newline
    Il sistema mostra l'errore e interrompe l'operazione.\newline
    Fine \\ \hline
    \textbf{Scenario alternativo} & 
    6a. Precondizione: campi vuoti o composti da soli spazi.\newline
    Il sistema mostra l'errore e interrompe l'operazione.\newline
    Fine \\ \hline
    \textbf{Scenario alternativo} & 
    7a. Precondizione: username già in uso (globale).\newline
    Il sistema mostra l'errore e interrompe l'operazione.\newline
    Fine \\ \hline 
    \textbf{Scenario alternativo} & 
    9a. Precondizione: Errore di salvataggio (I/O).\newline
    Il sistema rimuove il profilo dalla memoria (rollback), mostra l'errore e interrompe l'operazione.\newline
    Fine \\ \hline
\end{tabularx}

\renewcommand{\arraystretch}{1.4} 

\begin{tabularx}{\textwidth}{|l|X|}
    \hline
    \textbf{Nome}               & \textbf{UC22 - Login Fruitore} \\ \hline
    \textbf{Attore}             & Fruitore \\ \hline
    \endfirsthead 
    \hline
    \multicolumn{2}{|l|}{\textit{Continua dalla pagina precedente...}} \\ \hline
    \endhead 
    \hline
    \multicolumn{2}{|r|}{\textit{Continua nella pagina successiva...}} \\ \hline
    \endfoot 
    \hline
    \endlastfoot
    \textbf{Scenario principale} & 
    Precondizione: Il fruitore possiede un account e non è autenticato.\newline
    Trigger: Selezione "Login (account esistente)".\newline
    1. Il sistema mostra la schermata di login.\newline
    2. Il fruitore inserisce username e password.\newline
    3. Il sistema verifica la corrispondenza delle credenziali.\newline
    4. Il sistema autentica il fruitore.\newline
    5. Il sistema conteggia le notifiche presenti nello spazio personale e mostra il benvenuto con il relativo avviso.\newline
    Postcondizione: Il fruitore è autenticato nel sistema.\newline
    Fine \\ \hline
    \textbf{Scenario alternativo} & 
    2a. Precondizione: il fruitore digita "annulla" nel campo username.\newline
    Il sistema annulla l'operazione e torna al menu.\newline
    Fine \\ \hline
    \textbf{Scenario alternativo} & 
    3a. Precondizione: le credenziali sono errate (incluso inserire "annulla" come password).\newline
    Il sistema mostra l'errore e i tentativi rimasti.\newline
    Torna al punto 2 (se tentativi $<$ 3), altrimenti il sistema blocca l'accesso.\newline
    Fine \\ \hline
\end{tabularx}

\renewcommand{\arraystretch}{1.4} 

\begin{tabularx}{\textwidth}{|l|X|}
    \hline
    \textbf{Nome}               & \textbf{UC23 - Visualizza Bacheca (Fruitore)} \\ \hline
    \textbf{Attore}             & Fruitore \\ \hline
    \endfirsthead 
    \hline
    \multicolumn{2}{|l|}{\textit{Continua dalla pagina precedente...}} \\ \hline
    \endhead 
    \hline
    \multicolumn{2}{|r|}{\textit{Continua nella pagina successiva...}} \\ \hline
    \endfoot 
    \hline
    \endlastfoot
    \textbf{Scenario principale} & 
    Precondizione: Il fruitore è autenticato.\newline
    Trigger: Selezione "Visualizza bacheca".\newline
    1. Il fruitore richiede la visualizzazione.\newline
    2. Il sistema recupera tutte le proposte in stato APERTA.\newline
    3. Il sistema raggruppa le proposte per categoria.\newline
    4. Il sistema mostra dettagli, posti occupati, scadenze e tutti i campi informativi per ogni proposta.\newline
    5. Il sistema evidenzia con l'etichetta \textcolor{red}{[ISCRITTO]} le iniziative a cui l'utente loggato ha già aderito.\newline
    6. Il fruitore preme INVIO per terminare la consultazione.\newline
    Postcondizione: Nessuna modifica allo stato del sistema.\newline
    Fine \\ \hline  
    \textbf{Scenario alternativo} & 
    2a. Precondizione: nessuna proposta APERTA trovata nell'archivio.\newline
    Il sistema comunica "(Nessuna proposta aperta in bacheca)".\newline
    Fine \\ \hline
\end{tabularx}

\renewcommand{\arraystretch}{1.4} 

\begin{tabularx}{\textwidth}{|l|X|}
    \hline
    \textbf{Nome}               & \textbf{UC24 - Iscrizione a proposta} \\ \hline
    \textbf{Attore}             & Fruitore \\ \hline
    \endfirsthead 
    \hline
    \multicolumn{2}{|l|}{\textit{Continua dalla pagina precedente...}} \\ \hline
    \endhead 
    \hline
    \multicolumn{2}{|r|}{\textit{Continua nella pagina successiva...}} \\ \hline
    \endfoot 
    \hline
    \endlastfoot
    \textbf{Scenario principale} & 
    Precondizione: Il fruitore è autenticato.\newline
    Trigger: Selezione "Aderisci a una proposta".\newline
    1. Il fruitore avvia la procedura di iscrizione.\newline
    2. Il sistema estrae le proposte APERTE dall'archivio.\newline
    3. Il sistema verifica che ci sia almeno una proposta APERTA nella bacheca globale.\newline
    4. Il sistema filtra le proposte escludendo quelle a cui l'utente è già iscritto.\newline
    5. Il sistema verifica che la lista filtrata non sia vuota.\newline
    6. Il fruitore inserisce l'ID della proposta scelta.\newline
    7. Il sistema verifica che l'input sia numerico e diverso da "0".\newline
    8. Il sistema individua la proposta e verifica che sia in stato APERTA.\newline
    9. Il sistema verifica che la data odierna sia $\le$ al "Termine ultimo di iscrizione".\newline
    10. Il sistema verifica che ci siano posti disponibili.\newline
    11. Il sistema aggiunge lo username del fruitore alla lista degli aderenti.\newline
    12. Il sistema salva i dati in modo persistente.\newline
    13. Il sistema conferma l'avvenuta iscrizione all'utente.\newline
    Postcondizione: L'utente è iscritto e i posti disponibili scalano.\newline
    Fine \\ \hline
    \textbf{Scenario alternativo} & 
    3a. Precondizione: l'archivio globale non contiene alcuna proposta in stato APERTA.\newline
    Il sistema lo comunica all'utente e termina l'operazione.\newline
    Fine \\ \hline
    \textbf{Scenario alternativo} & 
    5a. Precondizione: il fruitore è già iscritto a tutte le proposte APERTE (lista filtrata vuota).\newline
    Il sistema lo comunica all'utente e termina l'operazione.\newline
    Fine \\ \hline
    \textbf{Scenario alternativo} & 
    7a. Precondizione: il fruitore digita "0" per annullare l'inserimento.\newline
    Il sistema annulla l'operazione.\newline
    Fine \\ \hline
    \textbf{Scenario alternativo} & 
    7b. Precondizione: l'input inserito non è un numero intero.\newline
    Il sistema mostra il messaggio "ID non valido".\newline
    Fine \\ \hline
    \textbf{Scenario alternativo} & 
    8a. Precondizione: l'ID non corrisponde a nessuna proposta o la proposta non è più APERTA.\newline
    Il sistema mostra un messaggio di errore.\newline
    Fine \\ \hline
    \textbf{Scenario alternativo} & 
    9a. Precondizione: il termine di iscrizione è scaduto rispetto alla data odierna.\newline
    Il sistema blocca l'iscrizione e mostra un messaggio di errore.\newline
    Fine \\ \hline
    \textbf{Scenario alternativo} & 
    10a. Precondizione: i posti sono esauriti (o il campo numero partecipanti è corrotto).\newline
    Il sistema blocca l'iscrizione e mostra un messaggio di errore.\newline
    Fine \\ \hline
    \textbf{Scenario alternativo} & 
    12a. Precondizione: Errore di salvataggio (I/O).\newline
    Il sistema annulla l'iscrizione in memoria (rollback) e mostra un messaggio di errore.\newline
    Fine \\ \hline
\end{tabularx}

\renewcommand{\arraystretch}{1.4} 

\begin{tabularx}{\textwidth}{|l|X|}
    \hline
    \textbf{Nome}               & \textbf{UC25 - Gestione Spazio Personale (Notifiche)} \\ \hline
    \textbf{Attore}             & Fruitore \\ \hline
    \endfirsthead 
    \hline
    \multicolumn{2}{|l|}{\textit{Continua dalla pagina precedente...}} \\ \hline
    \endhead 
    \hline
    \multicolumn{2}{|r|}{\textit{Continua nella pagina successiva...}} \\ \hline
    \endfoot 
    \hline
    \endlastfoot
    \textbf{Scenario principale} & 
    Precondizione: Il fruitore è autenticato.\newline
    Trigger: Selezione "Spazio personale".\newline
    1. Il fruitore accede allo spazio personale.\newline
    2. Il sistema mostra l'elenco delle notifiche e le opzioni del menu ("c", "t", "0").\newline
    3. Il fruitore inserisce l'opzione "c" (Cancella singola notifica).\newline
    4. Il sistema verifica che la lista delle notifiche non sia vuota.\newline
    5. Il sistema richiede l'inserimento dell'ID della notifica.\newline
    6. Il fruitore inserisce l'ID.\newline
    7. Il sistema individua e rimuove la notifica dalla memoria.\newline
    8. Il sistema salva i dati in modo persistente.\newline
    9. Il sistema conferma la cancellazione all'utente.\newline
    Torna al punto 2 \\ \hline
    \textbf{Scenario alternativo} & 
    2a. Precondizione: la lista delle notifiche dell'utente è vuota.\newline
    Il sistema stampa "(Nessuna notifica)" ma mostra comunque il menu operativo.\newline
    Prosegue al punto 3 \\ \hline
    \textbf{Scenario alternativo} & 
    3a. Precondizione: il fruitore seleziona "t" (Cancellazione totale).\newline
    Il sistema verifica che la lista non sia vuota, chiede conferma (s/n). L'utente risponde "s". Il sistema svuota la lista in memoria.\newline
    Salta al punto 8 \\ \hline
    \textbf{Scenario alternativo} & 
    3b. Precondizione: il fruitore seleziona "t", ma annulla la conferma (risponde "n").\newline
    Il sistema annulla l'operazione.\newline
    Torna al punto 2 \\ \hline
    \textbf{Scenario alternativo} & 
    3c. Precondizione: il fruitore seleziona "0. Torna al menu principale".\newline
    Il sistema chiude lo spazio personale e torna al menu precedente.\newline
    Fine \\ \hline
    \textbf{Scenario alternativo} & 
    3d. Precondizione: il fruitore inserisce un'opzione non valida (es. lettera errata).\newline
    Il sistema mostra l'errore "Scelta non valida".\newline
    Torna al punto 2 \\ \hline
    \textbf{Scenario alternativo} & 
    4a. Precondizione: il fruitore seleziona "c" o "t", ma la lista delle notifiche è vuota.\newline
    Il sistema mostra l'errore "Nessuna notifica da cancellare".\newline
    Torna al punto 2 \\ \hline
    \textbf{Scenario alternativo} & 
    7a. Precondizione: l'ID inserito non è numerico o non corrisponde ad alcuna notifica in memoria.\newline
    Il sistema mostra un errore.\newline
    Torna al punto 2 \\ \hline
    \textbf{Scenario alternativo} & 
    8a. Precondizione: Errore di salvataggio (I/O).\newline
    Il sistema ripristina in memoria la lista originale intatta (rollback) e mostra un errore.\newline
    Torna al punto 2 \\ \hline
\end{tabularx}

\renewcommand{\arraystretch}{1.4} 

\begin{tabularx}{\textwidth}{|l|X|}
    \hline
    \textbf{Nome}               & \textbf{UC26 - Visualizza Archivio Proposte} \\ \hline
    \textbf{Attore}             & Configuratore \\ \hline
    \endfirsthead 
    \hline
    \multicolumn{2}{|l|}{\textit{Continua dalla pagina precedente...}} \\ \hline
    \endhead 
    \hline
    \multicolumn{2}{|r|}{\textit{Continua nella pagina successiva...}} \\ \hline
    \endfoot 
    \hline
    \endlastfoot
    \textbf{Scenario principale} & 
    Precondizione: Il configuratore è autenticato.\newline
    Trigger: Selezione "Visualizza archivio proposte".\newline
    1. Il configuratore avvia la richiesta di visualizzazione.\newline
    2. Il sistema recupera l'intero archivio (escludendo gli stati interni BOZZA e VALIDA).\newline
    3. Il sistema raggruppa dinamicamente le proposte in base allo stato (APERTA, CONFERMATA, ANNULLATA, CONCLUSA, RITIRATA).\newline
    4. Il sistema mostra i dettagli anagrafici, i posti occupati, la lista degli aderenti e lo storico completo dei passaggi di stato per ogni proposta.\newline
    5. Il configuratore preme INVIO per terminare.\newline
    Postcondizione: Nessuna modifica allo stato del sistema.\newline
    Fine \\ \hline 
    \textbf{Scenario alternativo} & 
    2a. Precondizione: l'archivio è vuoto.\newline
    Il sistema segnala "(Nessuna proposta nell'archivio)" e termina l'operazione.\newline
    Fine \\ \hline
\end{tabularx}

\begin{tabularx}{\textwidth}{|l|X|}
    \hline
    \textbf{Nome}               & \textbf{UC27 - Ritira proposta in archivio} \\ \hline
    \textbf{Attore}             & Configuratore \\ \hline
    \endfirsthead 

    \hline
    \multicolumn{2}{|l|}{\textit{Continua dalla pagina precedente...}} \\ \hline
    \endhead 

    \hline
    \multicolumn{2}{|r|}{\textit{Continua nella pagina successiva...}} \\ \hline
    \endfoot 

    \hline
    \endlastfoot

    \textbf{Scenario principale} & 
    Precondizione: Il configuratore è autenticato. L'archivio contiene almeno una proposta in stato APERTA o CONFERMATA.\newline
    Trigger: Selezione dell'opzione "Ritira proposta".\newline
    1. Il sistema mostra le proposte in stato APERTA o CONFERMATA presenti nell'archivio.\newline
    2. Il configuratore inserisce l'ID della proposta da ritirare.\newline
    3. Il sistema estrae la proposta e verifica che il suo stato attuale consenta il ritiro.\newline
    4. Il sistema verifica che la data odierna sia strettamente precedente alla "Data inizio" dell'evento.\newline
    5. Il sistema transita lo stato della proposta in RITIRATA e registra il passaggio nello storico.\newline
    6. Il sistema invia automaticamente una notifica di ritiro nello spazio personale di tutti i fruitori iscritti (aderenti).\newline
    7. Il sistema salva le modifiche in modo persistente.\newline
    8. Il sistema mostra la conferma dell'operazione. \newline
    Fine \\ \hline
    
    \textbf{Scenario alternativo} & 
    2a. Precondizione: l'ID inserito non corrisponde ad alcuna proposta o la proposta non è nell'archivio.\newline
    Il sistema mostra un messaggio di errore e interrompe l'operazione. \newline
    Fine \\ \hline

    \textbf{Scenario alternativo} & 
    3a. Precondizione: la proposta si trova in uno stato diverso da APERTA o CONFERMATA (es. già CONCLUSA o ANNULLATA).\newline
    Il sistema mostra l'errore "Il ritiro è consentito solo per proposte APERTE o CONFERMATE" e interrompe l'operazione. \newline
    Fine \\ \hline

    \textbf{Scenario alternativo} & 
    4a. Precondizione: la "Data inizio" dell'evento è già trascorsa o corrisponde alla data odierna (oggi $\ge$ Data inizio).\newline
    Il sistema blocca il ritiro indicando che l'evento è imminente o passato. \newline
    Fine \\ \hline

    \textbf{Scenario alternativo} & 
    7a. Precondizione: Errore di salvataggio (I/O).\newline
    Il sistema esegue un rollback ricaricando i dati precedenti dal disco per ripristinare sia lo stato della proposta che le notifiche, poi mostra un errore. \newline
    Fine \\ \hline
\end{tabularx}

\renewcommand{\arraystretch}{1.4} 

\begin{tabularx}{\textwidth}{|l|X|}
    \hline
    \textbf{Nome}               & \textbf{UC28 - Disdici iscrizione a proposta} \\ \hline
    \textbf{Attore}             & Fruitore \\ \hline
    \endfirsthead 

    \hline
    \multicolumn{2}{|l|}{\textit{Continua dalla pagina precedente...}} \\ \hline
    \endhead 

    \hline
    \multicolumn{2}{|r|}{\textit{Continua nella pagina successiva...}} \\ \hline
    \endfoot 

    \hline
    \endlastfoot

    \textbf{Scenario principale} & 
    Precondizione: Il fruitore è autenticato e risulta iscritto ad almeno una proposta.\newline
    Trigger: Selezione dell'opzione per disdire un'iscrizione.\newline
    1. Il sistema mostra l'elenco delle proposte APERTE a cui l'utente è iscritto.\newline
    2. Il fruitore inserisce l'ID della proposta da cui desidera disiscriversi.\newline
    3. Il sistema verifica che la proposta sia in stato APERTA.\newline
    4. Il sistema verifica che la data odierna sia minore o uguale al "Termine ultimo di iscrizione".\newline
    5. Il sistema rimuove lo username del fruitore dalla lista degli aderenti della proposta.\newline
    6. Il sistema salva le modifiche in modo persistente.\newline
    7. Il sistema mostra la conferma dell'avvenuta disdetta. \newline
    Fine \\ \hline
    
    \textbf{Scenario alternativo} & 
    2a. Precondizione: l'ID non è valido oppure il fruitore non risulta iscritto alla proposta indicata.\newline
    Il sistema mostra un errore. \newline
    Fine \\ \hline

    \textbf{Scenario alternativo} & 
    3a. Precondizione: la proposta ha cambiato stato (non è più APERTA).\newline
    Il sistema mostra un errore indicando l'impossibilità di operare. \newline
    Fine \\ \hline

    \textbf{Scenario alternativo} & 
    4a. Precondizione: il "Termine ultimo di iscrizione" è scaduto rispetto alla data odierna.\newline
    Il sistema blocca la disdetta informando che il termine è oltrepassato. \newline
    Fine \\ \hline

    \textbf{Scenario alternativo} & 
    6a. Precondizione: Errore di salvataggio (I/O).\newline
    Il sistema esegue un rollback reinserendo il fruitore tra gli aderenti in memoria e mostra un messaggio di errore. \newline
    Fine \\ \hline
\end{tabularx}

\begin{tabularx}{\textwidth}{|l|X|}
    \hline
    \textbf{Nome} & \textbf{UC29 – Importa singolo file batch} \\ \hline
    \textbf{Attore} & Configuratore \\ \hline
    \endfirsthead 

    \hline
    \multicolumn{2}{|l|}{\textit{Continua dalla pagina precedente...}} \\ \hline
    \endhead 

    \hline
    \multicolumn{2}{|r|}{\textit{Continua nella pagina successiva...}} \\ \hline
    \endfoot 

    \hline
    \endlastfoot

    \textbf{Scenario principale} & 
    Precondizione: Il configuratore è autenticato. Il file batch esiste nel filesystem ed è leggibile.\newline
    Trigger: Selezione dell'opzione "Importa da file batch" $\rightarrow$ "Importa singolo file".\newline
    1. Il sistema mostra la descrizione del formato batch supportato (CAMPO\_COMUNE, CATEGORIA, PROPOSTA).\newline
    2. Il configuratore inserisce il percorso del file (assoluto o relativo alla working directory).\newline
    3. Il sistema legge tutte le righe del file.\newline
    4. Il sistema scorre le righe, ignorando quelle vuote e quelle che iniziano con \#.\newline
    5. Per ogni riga significativa, il sistema identifica il comando (primo token prima di |).\newline
    6. Per ogni riga CAMPO\_COMUNE: il sistema crea il campo comune con nome e obbligatorietà indicati e salva su disco.\newline
    7. Per ogni riga CATEGORIA: il sistema crea la categoria; se la riga contiene token CAMPO\_SPECIFICO, aggiunge i campi specifici inline prima del salvataggio, poi salva su disco.\newline
    8. Per ogni riga PROPOSTA: il sistema costruisce lo snapshot dei campi (BASE $\rightarrow$ COMUNI $\rightarrow$ SPECIFICI della categoria), assegna i valori indicati con il formato campo = valore, valida la proposta con le stesse regole della modalità interattiva e, se valida, la pubblica direttamente in bacheca salvando su disco.\newline
    9. Al termine, il sistema raccoglie il resoconto in un \textit{BatchRisultato} (numero righe elaborate, successi, warning, errori).\newline
    10. Il sistema mostra il resoconto completo all'utente.\newline
    Fine \\ \hline
    
    \textbf{Scenario alternativo} & 
    2a. Precondizione: il percorso è vuoto o l'utente digita "annulla".\newline
    Il sistema annulla l'operazione senza modifiche.\newline
    Fine \\ \hline

    \textbf{Scenario alternativo} & 
    3a. Precondizione: il file non esiste o non è accessibile (IOException).\newline
    Il sistema mostra un messaggio di errore con indicazione del percorso e interrompe l'operazione.\newline
    Fine \\ \hline

    \textbf{Scenario alternativo} & 
    5a. Precondizione: il comando della riga non è riconosciuto.\newline
    Il sistema registra un warning con numero di riga e prosegue con la riga successiva.\newline
    Fine (riga) \\ \hline

    \textbf{Scenario alternativo} & 
    6a. Precondizione: il campo comune ha nome vuoto, oppure esiste già un campo con lo stesso nome.\newline
    Il sistema registra un warning e salta la riga.\newline
    Fine (riga) \\ \hline

    \textbf{Scenario alternativo} & 
    6b. Precondizione: errore di salvataggio (I/O) dopo l'aggiunta in memoria.\newline
    Il sistema esegue il rollback in memoria, registra un errore e prosegue.\newline
    Fine (riga) \\ \hline

    \textbf{Scenario alternativo} & 
    7a. Precondizione: la categoria esiste già.\newline
    Il sistema registra un warning e salta la riga.\newline
    Fine (riga) \\ \hline

    \textbf{Scenario alternativo} & 
    7b. Precondizione: un campo specifico inline non è valido.\newline
    Il sistema registra un warning per quel campo specifico e lo salta; la categoria viene creata senza di esso.\newline
    Fine (riga) \\ \hline

    \textbf{Scenario alternativo} & 
    7c. Precondizione: errore di salvataggio (I/O) della categoria.\newline
    Il sistema esegue il rollback in memoria, registra un errore e prosegue.\newline
    Fine (riga) \\ \hline

    \textbf{Scenario alternativo} & 
    8a. Precondizione: la categoria indicata nella riga PROPOSTA non esiste.\newline
    Il sistema registra un errore e salta la riga.\newline
    Fine (riga) \\ \hline

    \textbf{Scenario alternativo} & 
    8b. Precondizione: token campo = valore non validi o campo non esistente.\newline
    Il sistema registra un warning per ogni token non riconosciuto e prosegue con i restanti campi.\newline
    Fine (riga) \\ \hline

    \textbf{Scenario alternativo} & 
    8c. Precondizione: la proposta risulta non valida dopo la validazione.\newline
    Il sistema registra un errore con i problemi rilevati. La proposta non viene pubblicata.\newline
    Fine (riga) \\ \hline

    \textbf{Scenario alternativo} & 
    8d. Precondizione: errore di salvataggio (I/O) della proposta pubblicata.\newline
    Il sistema esegue il rollback in memoria, registra un errore e prosegue.\newline
    Fine (riga) \\ \hline
\end{tabularx}

\renewcommand{\arraystretch}{1.4} 

\begin{tabularx}{\textwidth}{|l|X|}
    \hline
    \textbf{Nome} & \textbf{UC30 – Importa più file batch in sequenza} \\ \hline
    \textbf{Attore} & Configuratore \\ \hline
    \endfirsthead 

    \hline
    \multicolumn{2}{|l|}{\textit{Continua dalla pagina precedente...}} \\ \hline
    \endhead 

    \hline
    \multicolumn{2}{|r|}{\textit{Continua nella pagina successiva...}} \\ \hline
    \endfoot 

    \hline
    \endlastfoot

    \textbf{Scenario principale} & 
    Precondizione: Il configuratore è autenticato.\newline
    Trigger: Selezione dell'opzione "Importa da file batch" $\rightarrow$ "Importa più file in sequenza".\newline
    1. Il sistema chiede di inserire i percorsi dei file batch uno alla volta.\newline
    2. Il configuratore inserisce i percorsi (uno per invio).\newline
    3. Il configuratore preme INVIO con il campo vuoto per confermare la lista.\newline
    4. Il sistema mostra il numero di file selezionati e avvia l'elaborazione sequenziale.\newline
    5. Per ciascun file della lista, il sistema esegue l'elaborazione descritta in UC29 (passi 3–9).\newline
    6. I risultati parziali di ogni file vengono aggregati in un unico \textit{BatchRisultato} complessivo tramite il metodo \textit{aggiungi(parziale)}.\newline
    7. Al termine di tutti i file, il sistema mostra il resoconto aggregato con il numero totale di righe elaborate, successi, warning ed errori, e l'elenco dei file processati.\newline
    Fine \\ \hline
    
    \textbf{Scenario alternativo} & 
    2a. Precondizione: l'utente digita "annulla" invece di un percorso.\newline
    Il sistema annulla l'intera operazione senza apportare modifiche.\newline
    Fine \\ \hline

    \textbf{Scenario alternativo} & 
    3a. Precondizione: l'utente preme INVIO senza aver inserito alcun percorso.\newline
    Il sistema mostra un errore "Nessun file specificato" e interrompe l'operazione.\newline
    Fine \\ \hline

    \textbf{Scenario alternativo} & 
    5a. Precondizione: uno dei file indicati non esiste o non è leggibile (IOException).\newline
    Il sistema registra un errore nel resoconto aggregato (indicando il file) e continua con il file successivo nella lista.\newline
    Fine (file) \\ \hline

    \textbf{Scenario alternativo} & 
    5b. Precondizione: errori su singole righe di un file (comandi errati, duplicati, validazione fallita).\newline
    Si applicano le stesse alternative di UC29 (5a–8d) con i relativi rollback per riga. L'elaborazione prosegue per le righe e i file successivi.\newline
    Fine \\ \hline
\end{tabularx}
\clearpage
\section{Casi d'uso grafici}
\begin{figure}[htbp]
    \centering
    \makebox[\textwidth][c]{
    \centering
        \includegraphics[width=1.3\textwidth]{immagini/casi uso finale 1.png} 
    }
    \label{fig:usecase_vfinal1}
\end{figure}

\begin{figure}[htbp]
    \centering
    \makebox[\textwidth][c]{
        \includegraphics[width=1.3\textwidth]{immagini/casi uso finale 2.png} 
    }
    \label{fig:usecase_vfinal2}
\end{figure}

\begin{figure}[htbp]
    \centering
    \makebox[\textwidth][c]{
        \includegraphics[width=1.3\textwidth]{immagini/casi uso finale 3.png} 
    }
    \label{fig:usecase_vfinal3}
\end{figure}
\clearpage
\section{Diagramma delle classi}

\begin{figure}[p]
    \centering
    \includegraphics[width=\textwidth]{immagini/uml finale model.png}
    \label{fig:diagramma_classi_vfinal_model}
\end{figure}

\begin{figure}[p]
    \centering
    \includegraphics[width=0.9\textwidth]{immagini/uml finale view controller persistence.png}
    \label{fig:diagramma_classi_vfinal_view_controller_persistence}
\end{figure}
\clearpage
\section{Manuale di installazione e uso}

\end{document}